\section{Making a graph law reflection positive, and replacing a law at fixed marginals}\label{sec:completion-repair}

\paragraph{What this section proves and how to read it.}
A graph law need not be reflection positive: its cut operator can have a
negative part of size about $gp^{-1/2}$. This section proves three
constructions that, with the spectral caps, produce the objects
Theorem~C needs. A \emph{completion} adds four smaller summand
types (folds, matchings, merges, constants, with weights in powers
of $r=p^{-1/2}$) to make the law reflection positive at every cut.
\emph{Polarization} dominates two summands by two others in the
$2\times2$ block order. A \emph{replacement at fixed marginals}
(called a \emph{repair} in Part~V) reweights a coupling by
$1+f(x)+g(y)$ to prescribed one-half marginals with the same
support; ``repair'' never means a small perturbation, and a
repaired record keeps the actual half marginals of the record it
replaces, so differences annihilate constants. Inputs: spectral
caps (\cref{sec:spectral}). Outputs: the completion, polarization
and repair theorems applied in Part~V. A first reading should
follow the chamber geometry and the three statements below; weight
calculations are reference material.

\begin{enumerate}[label=(\arabic*)]
\item \emph{The five-type completion} (\cref{lem:five-type}). Adding to a
 root law its folds, matchings, merges and constants, with weights in powers
 of $r=p^{-1/2}$ fixed by four scalar inequalities, gives a law that is
 reflection positive at every chamber cut; adding one-layer diagonals
 restores the reserved cut $\iota$. This is (CP1) of \cref{def:cp1} for the
 middle atoms (\cref{lem:front12:reference}), the upper atoms
 (\cref{lem:upper-reference-rp}) and the tail atoms (\cref{sec:tail-index}).
 Its joint form (f) completes rows whose fold and matching kernels are bounded
 together, in one Gram matrix, with a mean allowance $t_Y>0$ for each row
 subject to $t_Y\le H^{-1}$ and $M\ge H+t_Y$
 (\cref{cor:joint-completion-blocks}); the bands of the upper range use it.
\item \emph{Polarization} (\cref{lem:tr:polarization}). At the chamber cut
 $\vartheta_0$ the two nonaligned fold records are dominated by the two
 matching records plus a merge correction of one more order in $r$. This is
 the first-stage order (CP2)(a) of \cref{def:cp2} (\cref{sec:middle-frame}),
 the matching bridge of the variance credits (\cref{sec:variance}) and, in all
 degrees, the fold coupon of the mean credit (\cref{prop:record:mean-credit}).
\item \emph{Additive marginal repair} (\cref{lem:tr:repair}) with the
 centred norm of an additive reweighting
 (\cref{lem:record:additive-cancellation}). A coupling with the required support whose prescribed target
 marginals are $\mu_X=(1+u)\nu_X$ and $\mu_Y=(1+v)\nu_Y$ with
 $|u|\le U_X$, $|v|\le U_Y$ (and the $L^2$ bounds of \cref{lem:tr:repair})
 is reweighted by $1+f(x)+g(y)$ to a coupling with exactly those marginals
 and the same support, and its centred conditional norm stays of the same order; the
 centred correction is bounded in the actual marginal space, with separate
 symmetric and skew parts and, against a law with the same marginals, without a
 marginal-variance term (\cref{lem:record:centered-correction}). If a coupling
 with the prescribed marginals and the required support is at hand, a positive
 kernel is completed by mixing it in (\cref{cor:markov-completion}). Apart from
 the matching records of the middle range, whose law keeps the half marginals of
 the records it replaces without a repair (\cref{lem:record:e-law}), the
 comparison records of (CP4) of \cref{def:cp4} are such repaired records or
 finite combinations of them: the signed fold law of the middle range combines
 three positive components completed through an accepting coupling
 (\cref{sec:middle-records}); the middle merge repairs use \cref{lem:tr:repair}
 (\cref{lem:merge-repair}); the upper range uses one-way repairs (imposing edges in one
direction only across the reserved cut), dense
 comparisons repaired by the initial-segment bound of \cref{lem:tr:repair}(g)
 and a convex mixture of such repairs (\cref{sec:upper-comparison}); and the
 tail uses one-way repairs
 (\cref{sec:tail-index}).
\end{enumerate}
A repaired record keeps the actual half marginals of the record it replaces,
and so does a signed combination of such records with total coefficient one;
so every record difference annihilates the constants on both sides of the
reserved cut. This is what makes the sharp prices, the variance credits and
the common mean of Part~V (the empty-support component retained
in \cref{lem:record:common-mean}(a)) possible. \Cref{subsec:chamber-geometry} fixes the
chamber geometry of the records, which all three statements use.

\subsection{How records sit inside chambers (supports and symmetries)}\label{subsec:chamber-geometry}

Throughout this section the patterns below act in one chamber
coordinate $\Z/6\Z$ (two in \cref{lem:five-type}(d)); all remaining slot data
of a configuration are carried by the objects in $\mathfrak X$. We use the
rotation $\rot$, the chamber reflections $\vartheta_i(j)=2i-1-j$ with plus
halves $H_i=\{i,i+1,i+2\}$ ($\vartheta_{i+3}=\vartheta_i$), and the group
$\mathcal G_6$ of \cref{def:slot-array,lem:slot-array}(i): the translations
$u\mapsto u+2k$ and the reflections $\vartheta_0,\vartheta_1,\vartheta_2$
($u\mapsto5-u$, $1-u$, $3-u$).

An \emph{object} carries the labels of a finite set $\mathfrak X$, and a
\emph{configuration} is $\Psi=(\Psi_u)_{u\in\Z/6\Z}\in\mathfrak X^{\Z/6\Z}$. For a
cut $\vartheta=\vartheta_i$ we use the cut matrix $C_\vartheta$, the order
$\succeq_\vartheta$ and the cut form of \cref{def:cc:order} on configurations, and
write the diagonal cut form at a real function $g$ on $\mathfrak X^{H_i}$ as
\[
 Q^{\vartheta}_\kappa(g)=\int g(\Psi|_{H_i})\,g\bigl((\Psi\circ\vartheta)|_{H_i}\bigr)
 \,d\kappa(\Psi)=\ip g{C_\vartheta(\kappa)g},\qquad Q_\kappa=Q^{\vartheta_0}_\kappa .
\]
For a law $\lambda$ on $\mathfrak X^{H_i}$ the diagonal copy
$\operatorname{Diag}_\vartheta(\lambda)$
(\cref{lem:commoncut}(i)) is fixed \emph{samplewise} by
$\vartheta$, meaning $\Psi\circ\vartheta=\Psi$ for every
configuration $\Psi$ in its support, and its cut form is
$\E_\lambda g^2$.

\begin{definition}[Patterns and types]\label{def:patterns}
For a map $\phi\colon\Z/6\Z\to O$ into a finite set of objects and a law
$\lambda$ on $\mathfrak X^O$, $\Pi_\phi(\lambda)$ is the law of the configuration
$\omega\circ\phi=(\omega_{\phi(u)})_u$, $\omega\sim\lambda$; we write $\Pi_\phi$
when $\lambda$ is clear. The \emph{patterns}, listed as $(\phi(0),\dots,\phi(5))$,
are:
\begin{itemize}
\item the identity of $\Z/6\Z$ (type $\rm L$, the \emph{root});
\item the three \emph{folds} $F_0=(0,1,2,2,1,0)$, $F_1=(1,0,0,1,2,2)$,
 $F_2=(2,2,1,0,0,1)$ (type $\rm D$), with $O=\{0,1,2\}$;
\item the two \emph{matchings} $E_0=(0,1,1,2,2,0)$ and $E_1=(0,0,1,1,2,2)$
 (type $\rm S$), constant on the blocks of the pair partitions
 $\mathcal P_e=\{\{0,5\},\{1,2\},\{3,4\}\}$ and
 $\mathcal P_o=\{\{0,1\},\{2,3\},\{4,5\}\}$, whose $k$th blocks are
 $E_0^{-1}(k)$ and $E_1^{-1}(k)$;
\item the six \emph{merges} $G_{sk}=\one[E_s(\cdot)=k]$, $s\in\{0,1\}$,
 $k\in\{0,1,2\}$ (type $\rm R$), with values in $\{0,1\}$:
 $G_{00}=(1,0,0,0,0,1)$, $G_{01}=(0,1,1,0,0,0)$, $G_{02}=(0,0,0,1,1,0)$,
 $G_{10}=(1,1,0,0,0,0)$, $G_{11}=(0,0,1,1,0,0)$, $G_{12}=(0,0,0,0,1,1)$; we
 write $G_e=G_{00}$ and $G_o=G_{11}$;
\item the constant map $0$, counted twice as $C_0$ and $C_1$ (type $\rm C$),
 so that $\pi_{\rm C}$ in \cref{def:five-type-datum} contains two copies of
 $\Pi_0(\nu)$ and contributes the term $2a_{\rm C}z^2$ in
 \cref{lem:five-type}(a).
\end{itemize}
The types $\rm L,D,S,R,C$ have the \emph{multiplicities}
$1,3,2,6,2$ and the \emph{formal degrees} $0,1,1,2,3$, where the
formal degree is the exponent of $r$ in the graded weight of
\cref{lem:five-type}.
\end{definition}

The geometry of these patterns is the following finite check. The root graph
of the middle reference, with one or two layers, is the \emph{rotation lift}
$G_{\rm L}$ below; with two layers its uniform law is the twelve-label root of
\cref{sec:middle-reference}, and on the three late middle pieces the root is a
positive reweighting of this law, with the same support and the same
symmetries (\cref{sec:positive-roots}).

\begin{table}[ht]
\centering\small
\begin{tabular}{@{}llll@{}}
\toprule
pattern & plus half $(\Psi_0,\Psi_1,\Psi_2)$ & reflected half $(\Psi_5,\Psi_4,\Psi_3)$ & position\\
\midrule
$F_0$ & $(\omega_0,\omega_1,\omega_2)$ & $(\omega_0,\omega_1,\omega_2)$ & diagonal\\
$F_1$ & $(\omega_1,\omega_0,\omega_0)$ & $(\omega_2,\omega_2,\omega_1)$ & $\mathcal U\to\mathcal V$\\
$F_2$ & $(\omega_2,\omega_2,\omega_1)$ & $(\omega_1,\omega_0,\omega_0)$ & $\mathcal V\to\mathcal U$\\
$E_0$ & $(\omega_0,\omega_1,\omega_1)$ & $(\omega_0,\omega_2,\omega_2)$ & $\mathcal U\to\mathcal U$\\
$E_1$ & $(\omega_0,\omega_0,\omega_1)$ & $(\omega_2,\omega_2,\omega_1)$ & $\mathcal V\to\mathcal V$\\
$G_{00}=G_e$ & $(\omega_1,\omega_0,\omega_0)$ & $(\omega_1,\omega_0,\omega_0)$ & diagonal on $\mathcal U$\\
$G_{01}$ & $(\omega_0,\omega_1,\omega_1)$ & $(\omega_0,\omega_0,\omega_0)$ & $\mathcal U\to\mathcal Z$\\
$G_{02}$ & $(\omega_0,\omega_0,\omega_0)$ & $(\omega_0,\omega_1,\omega_1)$ & $\mathcal Z\to\mathcal U$\\
$G_{10}$ & $(\omega_1,\omega_1,\omega_0)$ & $(\omega_0,\omega_0,\omega_0)$ & $\mathcal V\to\mathcal Z$\\
$G_{11}=G_o$ & $(\omega_0,\omega_0,\omega_1)$ & $(\omega_0,\omega_0,\omega_1)$ & diagonal on $\mathcal V$\\
$G_{12}$ & $(\omega_0,\omega_0,\omega_0)$ & $(\omega_1,\omega_1,\omega_0)$ & $\mathcal Z\to\mathcal V$\\
$C_0,C_1$ & $(\omega_0,\omega_0,\omega_0)$ & $(\omega_0,\omega_0,\omega_0)$ & diagonal on $\mathcal Z$\\
\bottomrule
\end{tabular}
\caption{The halves of $\omega\circ\phi$ at $\vartheta_0$, for
$\omega=(\omega_0,\omega_1,\omega_2)$, and the corners
$\mathcal U=\{(b,y,y)\colon y\ne b\}$, $\mathcal V=\{(y,y,b)\colon y\ne b\}$,
$\mathcal Z=\{(b,b,b)\}$ of $\mathfrak X^{H_0}$ in which they lie when the
labels $\omega_o$ that occur are distinct. In the position column
$\mathcal A\to\mathcal B$ means that the plus half lies in $\mathcal A$ and
the reflected half in $\mathcal B$.}
\label{tab:corner}
\end{table}

\begin{table}[ht]
\centering\small
\begin{tabular}{@{}lccc@{}}
\toprule
$\alpha$ & $(F_0\circ\alpha,F_1\circ\alpha,F_2\circ\alpha)$ & $\psi_{0,\alpha}$ & $\psi_{1,\alpha}$\\
\midrule
$\vartheta_1\colon u\mapsto1-u$ & $(F_1,F_0,F_2)$ & $(0\ 1)$ & $(1\ 2)$\\
$\vartheta_2\colon u\mapsto3-u$ & $(F_2,F_1,F_0)$ & $(0\ 2)$ & $(0\ 1)$\\
$\vartheta_0\colon u\mapsto5-u$ & $(F_0,F_2,F_1)$ & $(1\ 2)$ & $(0\ 2)$\\
$u\mapsto u+2$ & $(F_2,F_0,F_1)$ & $(0\ 1\ 2)$ & $(0\ 1\ 2)$\\
\bottomrule
\end{tabular}
\caption{Precomposition with the generators of $\mathcal G_6$:
$E_s\circ\alpha=\psi_{s,\alpha}\circ E_s$, where $(0\ 1\ 2)$ is the cycle
$0\mapsto1\mapsto2\mapsto0$.}
\label{tab:reflection}
\end{table}

\begin{lemma}[Chamber geometry of the records]\label{lem:record:fold-identification}
Let $\Blk$ be a finite set of primes and $\Lambda$ a finite set of layers. Let
$G_{\rm L}$ be the oriented graph on $\Z/6\Z\times\Lambda$ with the edges
$(u,\ell)\to(\rot u,\ell')$ for all $u$ and all $\ell,\ell'\in\Lambda$ (all layer
lifts of the six rotation edges $0\to2\to4\to0$ and $1\to5\to3\to1$); let
$G_{\rm F}$ be the graph on $\{0,1,2\}\times\Lambda$ with the
edges $(0,\ell)\to(2,\ell')$ for all $\ell,\ell'\in\Lambda$,
and $G_{\rm E}$ the graph on $\{0,1,2\}\times\Lambda$
without edges. An object carries one label of $X_{\Blk}=\prod_{p\in\Blk}\Fp$ per
layer, so $\mathfrak X=X_{\Blk}^\Lambda$. The graph law $\lambda_{\Blk}(G_{\rm L})$ is a law on
$\mathfrak X^{\Z/6\Z}$; put $\lambda_{\rm F}=\lambda_{\Blk}(G_{\rm F})$ and
$\lambda_{\rm E}=\lambda_{\Blk}(G_{\rm E})$.
\begin{enumerate}[label=\textup{(\alph*)}]
\item The reflections $u\mapsto k-u$ $(k\in\{1,3,5\})$ and the translation
 $u\mapsto u+2$ map the edge set of $G_{\rm L}$ onto itself. Hence $\lambda_{\Blk}(G_{\rm L})$ is invariant under $\mathcal G_6$, and at each $\vartheta_i$ it is a coupling
 of two copies of one half law.
\item For $j\in\{0,1,2\}$, $F_j\circ\vartheta_{2j}=F_j$, and $F_j$ maps the
 points $2j,2j+1,2j+2$ of $H_{2j}$ (indices mod $6$) to $0,1,2$. Hence for every
 law $\lambda$ on $\mathfrak X^{\{0,1,2\}}$,
 $\Pi_{F_j}(\lambda)=\operatorname{Diag}_{\vartheta_{2j}}(\lambda^{(j)})$, where
 $\lambda^{(j)}$ is the law on $\mathfrak X^{H_{2j}}$ which puts the labels of the
 object $o$ at $2j+o$, $o\in\{0,1,2\}$. We call $\vartheta_{2j}$ the
 \emph{fold chamber} of $F_j$.
\item The subgraph of $G_{\rm L}$ induced on $H_{2j}\times\Lambda$, read in the
 order $(2j,2j+1,2j+2)$, is $G_{\rm F}$. Hence the half marginal of
 $\lambda_{\Blk}(G_{\rm L})$ on $H_{2j}$, read in this order, is a law on the support of
 $\lambda_{\rm F}$, and it does not depend on $j$.
\item Let $\Lambda=\{0,1\}$ and let $\iota$ be the layer swap, with plus half the
 layer $0$. If a law $\lambda$ on $\mathfrak X^O$ is invariant under the swap of
 the two layer labels of every object, then every image $\Pi_\phi(\lambda)$ is
 symmetric at $\iota$, and its actual half marginal at $\iota$ is the image
 under $\phi$ of the layer-zero marginal of $\lambda$. In particular
 $\Pi_{F_j}(\lambda_{\rm F})$ and $\Pi_{E_s}(\lambda_{\rm E})$ are symmetric at
 $\iota$; the actual half marginal of the first is $\Pi_{F_j}(m_0)$, where
 $m_0$ is the law of the layer-zero triple under $\lambda_{\rm F}$, that of the second the
 uniform law on the ordered triples of labels distinct at every prime of
 $\Blk$, placed by $E_s$.
\item Precomposition with $\alpha\in\mathcal G_6$ permutes the folds, maps
 $E_s$ to $\psi\circ E_s$ for a permutation $\psi=\psi_{s,\alpha}$ of
 $\{0,1,2\}$, maps $G_{sk}$ to $G_{s,\psi^{-1}(k)}$, and fixes the constant
 maps; for the generators the action is that of \cref{tab:reflection}. The
 fold $F_j$ and the merges $G_e$, $G_o$ are fixed by the following
 chamber reflections: $F_0$, $G_e$ and $G_o$ by $\vartheta_0$, $F_1$ by $\vartheta_2$,
 $F_2$ by $\vartheta_4=\vartheta_1$.
\item At $\vartheta_0$ the halves of $\omega\circ\phi$ are those of
 \cref{tab:corner}. The corners $\mathcal U$, $\mathcal V$, $\mathcal Z$ are
 pairwise disjoint.
\end{enumerate}
Parts (b), (e) and (f) use only the patterns, and hold for objects with labels
in any finite set $\mathfrak X$.
\end{lemma}

\begin{proof}
(a) Under $u\mapsto3-u$ the edges $0\to2$, $2\to4$, $4\to0$ go to $3\to1$,
$1\to5$, $5\to3$, and conversely; under $u\mapsto1-u$ they go to $1\to5$,
$5\to3$, $3\to1$, and under $u\mapsto5-u$ to $5\to3$, $3\to1$, $1\to5$; the
translation by two permutes each of the two directed triangles. Since all layer
lifts of every rotation edge are edges, these maps, acting on the first
factor, are automorphisms of $G_{\rm L}$, and they generate $\mathcal G_6$. A
graph automorphism permutes the injective strict-square labellings at every
prime, so it preserves $\lambda_{\Blk}(G_{\rm L})$. A law invariant under
$\vartheta_i$ has the same law on its two halves.

(b) At the points $2j,2j+1,2j+2$ the fold $F_j$ takes the values $0,1,2$
(for $F_2$ the points are $4,5,0$), and $\vartheta_{2j}(u)=4j-1-u$, that is
$5-u$, $3-u$ and $1-u$ for $j=0,1,2$; substitution in the tuples gives
$F_j(\vartheta_{2j}u)=F_j(u)$. So every configuration $\omega\circ F_j$ is fixed
by $\vartheta_{2j}$, and its plus half carries $\omega_o$ at $2j+o$.

(c) Inside $H_{2j}$ the only rotation edge is $2j\to2j+2$: $\rot(2j)=2j+2$,
while $\rot(2j+1)=2j-1$ and $\rot(2j+2)=2j+4$ lie outside $H_{2j}$, and the odd
position $2j+1$ is joined in $G_{\rm L}$ only to odd positions. With all layer lifts the induced
graph is $G_{\rm F}$ in the stated order. The restriction of an injective
strict-square labelling of $G_{\rm L}$ to $H_{2j}\times\Lambda$ is one of the
induced graph. The translation $u\mapsto u+2j$ maps $H_0$ onto $H_{2j}$ in order and
preserves $G_{\rm L}$, so by (a) the half marginal does not depend on $j$.

(d) The layer swap commutes with $\omega\mapsto\omega\circ\phi$, so the image of
a swap-invariant law is $\iota$-invariant, which is the symmetry of its cut
matrix at $\iota$; its plus half is the layer-zero part. The swap is an
automorphism of $G_{\rm F}$ (all layer lifts of $0\to2$ are edges) and of
$G_{\rm E}$, so $\lambda_{\rm F}$ and $\lambda_{\rm E}$ are swap invariant. Under
$\lambda_{\rm E}$ the six labels are uniform among those distinct at every
prime, so the layer-zero triple is uniform on the ordered triples of distinct
labels.

(e) Substitution in the tuples gives \cref{tab:reflection}; for instance
$E_0(1-u)=(1,0,0,2,2,1)$, which is $E_0$ with the names $0,1$ exchanged, and
$F_1(3-u)=F_1(u)$. Since the generators map each block of $\mathcal P_e$ and of
$\mathcal P_o$ onto a block of the same partition, so does every
$\alpha\in\mathcal G_6$, and $E_s\circ\alpha=\psi\circ E_s$ for the induced
permutation $\psi$ of the block names. Then
$G_{sk}(\alpha u)=\one[\psi(E_s(u))=k]=G_{s,\psi^{-1}(k)}(u)$. The last
sentence is (b) together with $G_{00}(5-u)=G_{00}(u)$ and
$G_{11}(5-u)=G_{11}(u)$.

(f) Substitution gives \cref{tab:corner}. If $(b,y,y)=(y',y',b')$ with
$y\ne b$ and $y'\ne b'$, then $y'=b$ and $y'=y$, contradicting $y\ne b$; a
point of $\mathcal Z$ has three equal entries, so in either form above it
has $y=b$ and lies in neither $\mathcal U$ nor $\mathcal V$. Hence the
corners are pairwise disjoint.
\end{proof}

The cut forms of all components at $\vartheta_0$ are read off
\cref{tab:corner}: the cut form of every component other than the root and the
aligned fold $F_0$ depends on a half-function $g$ only through its
restrictions to the corners $\mathcal U$, $\mathcal V$ and $\mathcal Z$.
In Part~V the two facts used most are (b), which says that the two-layer fold
record is the diagonal copy, at its own chamber, of its induced law, and (d),
which identifies its actual half marginal at the reserved cut.

\subsection{The five-type completion}\label{subsec:five-type}

A labelling law is made positive at a chamber cut in steps of one order of
$r=p^{-1/2}$ each. Relative to the diagonal of its own half law, its cut
operator is bounded below by $-\delta$ with $\delta$ of order $r$ (the
\emph{root envelope}, from the spectral caps). In the graph-law
applications that diagonal is supplied by the aligned fold
(\cref{lem:record:fold-identification}(b),(c)): a copy of the half law
placed symmetrically, which is fixed by the cut and positive there. In
\cref{lem:five-type} this is the general root-envelope hypothesis
$\Pi_{\rm L}\succeq_{\vartheta_0}-\delta\,\Pi_{F_0}$, with $\Pi_{F_0}$
built from $\pi_{\rm F}$. The aligned fold at one cut is a nonaligned
fold at the others, where its form couples the two corners
$\mathcal U,\mathcal V$; this coupling is dominated by the matchings up to the kernel
errors $e_{\rm F},e_{\rm E}$ of \cref{lem:five-type}(b); those errors, and
the mean-zero parts $U^\circ,V^\circ$ of (c), are dominated by the merges,
and the linear terms of the merges by the constants. So the
five types enter with the formal degrees $0,1,1,2,3$, and the constants may
carry degree three although the merges have degree two, because the last
condition $a_{\rm C}a_{\rm S}\ge a_{\rm R}^2$ of \cref{lem:five-type}(b)
has total degree $3+1=2+2$ in the graded weights.

\begin{definition}[Five-type datum]\label{def:five-type-datum}
A \emph{five-type datum} on a finite set $\mathfrak X$ consists of
\begin{enumerate}[label=\textup{(\roman*)}]
\item a probability measure $\nu$ on $\mathfrak X$, the law of the middle
 object $\omega_1$ in (iii) and (iv) below (the \emph{middle law}), and,
 for each $b\in\supp\nu$, a probability measure $\nu^b$ on $\mathfrak X$ (the
 \emph{conditional reference});
\item a symmetric probability measure $\Gamma$ on $\mathfrak X^2$ with
 $\Gamma(dx\,db)=\nu(db)\nu^b(dx)$;
\item an exchangeable probability measure $\pi_{\rm E}$ on
 $\mathfrak X^{\{0,1,2\}}$ with
 \[
  \pi_{\rm E}(d\omega_0\,d\omega_1\,d\omega_2)=\nu(d\omega_1)\,\nu^{\omega_1}(d\omega_0)\,
  \nu^{\omega_1}(d\omega_2)\,T_{\rm E}^{\omega_1}(\omega_0,\omega_2);
 \]
\item a probability measure $\pi_{\rm F}$ on $\mathfrak X^{\{0,1,2\}}$ with
 \[
  \pi_{\rm F}(d\omega_0\,d\omega_1\,d\omega_2)=\nu(d\omega_1)\,\nu^{\omega_1}(d\omega_0)\,
  \nu^{\omega_1}(d\omega_2)\,T_{\rm F}^{\omega_1}(\omega_0,\omega_2);
 \]
\item a probability measure $\pi_{\rm L}$ on $\mathfrak X^{\Z/6\Z}$ invariant
 under $\mathcal G_6$;
\end{enumerate}
with nonnegative kernels $T_{\rm E}^b,T_{\rm F}^b$. A kernel $T$ is the
operator $Th(x)=\int T(x,y)h(y)\,\nu^b(dy)$ on $L^2(\nu^b)$; norms, inner
products, means and the projection $\Jproj$ onto the constants refer to
$L^2(\nu^b)$. The \emph{components} of the datum are
$\Pi_{\rm L}=\pi_{\rm L}$, $\Pi_{F_j}=\Pi_{F_j}(\pi_{\rm F})$,
$\Pi_{E_s}=\Pi_{E_s}(\pi_{\rm E})$, $\Pi_{G_{sk}}=\Pi_{G_{sk}}(\Gamma)$ (the
objects $0,1$ of $\Gamma$) and $\Pi_{C_0}=\Pi_{C_1}=\Pi_0(\nu)$; $\pi_t$ is the
sum of the components of the type $t$, and for weights
$a=(a_{\rm L},a_{\rm D},a_{\rm S},a_{\rm R},a_{\rm C})\ge0$ the \emph{five-type
measure} is $\mathcal Q_a=\sum_ta_t\pi_t$.
\end{definition}

In the applications $\Gamma$ is the two-object marginal of $\pi_{\rm E}$ and
$\nu$, $\nu^b$ are its one-object marginal and conditional law (the merges and
constants use only some objects of the matching law); then (ii) holds, and
$T^b_{\rm E}$, the density of the conditional law of the objects $0,2$ given
$\omega_1=b$, is a coupling of $\nu^b$ with itself, $T_{\rm E}^b\one=
(T_{\rm E}^b)^*\one=\one$. The fold kernel $T_{\rm F}^b$ contains the density of
the middle law of $\pi_{\rm F}$ relative to $\nu$ and those of its conditional
outer marginals relative to $\nu^b$, so no equality of the middle laws of the
fold and matching laws is assumed.

\begin{definition}[Two coordinates]\label{def:two-coordinate-five-type}
Let $X$ be a finite set of labels, $\Lambda$ a finite set of layers and
$W=\Z/6\Z\times\Z/6\Z\times\Lambda$. Put $O_{\rm L}=\Z/6\Z$ and
$O_{\rm F}=O_{\rm E}=\{0,1,2\}$. The types $\rm L,D,S,R,C$ have
the \emph{field types} (the named graph $G_{\rm L},G_{\rm F}$ or
$G_{\rm E}$ from which the coordinate law is taken; ``field''
here does not refer to $\Fp$)
$\mathfrak f({\rm L})={\rm L}$, $\mathfrak f({\rm D})={\rm F}$,
$\mathfrak f({\rm S})=\mathfrak f({\rm R})=\mathfrak f({\rm C})={\rm E}$ and the
pattern sets $P_{\rm L}=\{{\rm id}\}$, $P_{\rm D}=\{F_0,F_1,F_2\}$,
$P_{\rm S}=\{E_0,E_1\}$, $P_{\rm R}=\{G_{sk}\}$, $P_{\rm C}=\{C_0,C_1\}$. Suppose
given, for each pair of field types, a probability law
$\lambda_{\mathfrak f_0\mathfrak f_1}$ on $X^{O_{\mathfrak f_0}\times O_{\mathfrak f_1}
\times\Lambda}$. A \emph{state} is $t=(t_0,t_1)$, of degree
$\deg t_0+\deg t_1$. For patterns $\phi=(\phi_0,\phi_1)\in P_{t_0}\times P_{t_1}$
the component $\Pi_{t,\phi}$ is the law of $e_\phi(y)$, where
$e_\phi(y)_{(u_0,u_1,\ell)}=y_{\phi_0(u_0),\phi_1(u_1),\ell}$ and
$y\sim\lambda_{\mathfrak f(t_0)\mathfrak f(t_1)}$; $\pi_t=\sum_\phi\Pi_{t,\phi}$, and
for weights $a_t\ge0$ we put $\mathcal Q=\sum_ta_t\pi_t$. The \emph{row data
of the coordinate $0$ at the field type $\mathfrak f'$} are the laws, on the
row space $\mathfrak X=X^{O_{\mathfrak f'}\times\Lambda}$, of the families of rows
$(y_{o,\cdot,\cdot})_o$ of $y\sim\lambda_{{\rm L}\mathfrak f'}$ (this is
$\pi_{\rm L}$), $y\sim\lambda_{{\rm F}\mathfrak f'}$ ($\pi_{\rm F}$) and
$y\sim\lambda_{{\rm E}\mathfrak f'}$ ($\pi_{\rm E}$), together with the law
$\Gamma$ of the rows $0,1$ and the law $\nu$ of the row $0$ of
$\lambda_{{\rm E}\mathfrak f'}$, and the conditional law $\nu^b$ of the row $0$
given that the row $1$ is $b$. The row data of the coordinate $1$ are defined
in the same way with the two coordinates exchanged. We write
$\vartheta_{c,i}$ for the reflection $\vartheta_i$ in the coordinate $c$.
\end{definition}

\begin{lemma}[Five-type completion]\label{lem:five-type}
Let $(\nu,\nu^b,\Gamma,\pi_{\rm E},\pi_{\rm F},\pi_{\rm L})$ be a five-type datum
on $\mathfrak X$, and $a\ge0$ weights.
\begin{enumerate}[label=\textup{(\alph*)}]
\item \emph{Invariance and cut forms.} $\mathcal Q_a$ is invariant under
 $\mathcal G_6$. For a real function $g$ on $\mathfrak X^{H_0}$ put
 $U_b(y)=g(b,y,y)$, $V_b(y)=g(y,y,b)$, $z_b=g(b,b,b)$, $u_b=\E_{\nu^b}U_b$ and
 $v_b=\E_{\nu^b}V_b$, and let $\E_b$ be the average over $b\sim\nu$. Then
 \begin{gather*}
  Q_{\Pi_{F_0}}(g)=\E_{\pi_{\rm F}}g(\omega)^2,\qquad
  Q_{\Pi_{F_1}}(g)=Q_{\Pi_{F_2}}(g)=\E_b\ip{U_b}{T_{\rm F}^bV_b},\\
  Q_{\Pi_{E_0}}(g)=\E_b\ip{U_b}{T_{\rm E}^bU_b},\qquad
  Q_{\Pi_{E_1}}(g)=\E_b\ip{V_b}{T_{\rm E}^bV_b},\\
  Q_{\Pi_{G_e}}(g)=\E_b\|U_b\|^2,\qquad Q_{\Pi_{G_o}}(g)=\E_b\|V_b\|^2,\\
  Q_{\Pi_{G_{01}}}(g)=Q_{\Pi_{G_{02}}}(g)=\E_b[z_bu_b],\qquad
  Q_{\Pi_{G_{10}}}(g)=Q_{\Pi_{G_{12}}}(g)=\E_b[z_bv_b],\qquad
  Q_{\Pi_{C_k}}(g)=\E_bz_b^2 .
 \end{gather*}
 For every $b\in\supp\nu$, every linear operator $T$ on
 $L^2(\nu^b)$, every $e\ge0$ with $\|T-\Jproj\|\le e$ and all
  real $\Phi,\Psi\in L^2(\nu^b)$,
 $\ip\Phi{T\Psi}\ge\E_{\nu^b}\Phi\,\E_{\nu^b}\Psi
 -e\|\Phi\|\|\Psi\|$.
\item \emph{Full-norm criterion.} Suppose that
 $\|T_{\rm F}^b-\Jproj\|\le e_{\rm F}$ and $\|T_{\rm E}^b-\Jproj\|\le e_{\rm E}$ for
 every $b\in\supp\nu$, and
 \[
  a_{\rm S}\ge a_{\rm D},\qquad a_{\rm R}\ge e_{\rm F}a_{\rm D}+e_{\rm E}a_{\rm S},
  \qquad a_{\rm C}a_{\rm S}\ge a_{\rm R}^2 .
 \]
 Then $\mathcal Q_a-a_{\rm L}\Pi_{\rm L}\succeq_{\vartheta_0}a_{\rm D}\Pi_{F_0}$. If
 moreover the root envelope $\Pi_{\rm L}\succeq_{\vartheta_0}-\delta\,\Pi_{F_0}$
 holds and $a_{\rm D}\ge\delta a_{\rm L}$ (no condition if $a_{\rm L}=0$), then
 $\mathcal Q_a$ is reflection
 positive at every chamber cut $\vartheta_i$, $i\in\Z/6\Z$. For $r>0$ and the
 \emph{graded weights} $(w_{\rm L},w_{\rm D}r,w_{\rm S}r,w_{\rm R}r^2,w_{\rm C}r^3)$,
 with $e_{\rm F}=fr$, $e_{\rm E}=hr$ and $\delta=dr$, the conditions read
 \[
  w_{\rm D}\ge d\,w_{\rm L},\qquad w_{\rm S}\ge w_{\rm D},\qquad
  w_{\rm R}\ge fw_{\rm D}+hw_{\rm S},\qquad w_{\rm C}w_{\rm S}\ge w_{\rm R}^2 .
 \]
\item \emph{Schur criterion.} Suppose that for every $b\in\supp\nu$ the
 matching kernel has exact marginals, $T_{\rm E}^b\one=(T_{\rm E}^b)^*\one=\one$,
 with $\|T_{\rm E}^b-\Jproj\|\le\rho_{\rm E}$, and that the fold kernel
 $K_b=T_{\rm F}^b$ has the \emph{fold data}
 \[
  m_b:=\ip\one{K_b\one}\in[H^{-1},H],\qquad
  \|(I-\Jproj)K_b\one\|,\ \|(I-\Jproj)K_b^*\one\|\le\beta,\qquad
  \|(I-\Jproj)K_b(I-\Jproj)\|\le q_\perp .
 \]
 Let $a_{\rm D}>0$, $M=a_{\rm S}/a_{\rm D}$, $s=a_{\rm R}/a_{\rm D}$,
 $k=a_{\rm C}/a_{\rm D}$, and let $t_Y>0$ (the \emph{mean allowance})
 satisfy
 \[
  M\ge H+t_Y,\qquad t_Y\le H^{-1},\qquad
  s\ge q_\perp+M\rho_{\rm E}+\beta^2/t_Y,\qquad kM\ge s^2 .
 \]
 Then $\mathcal Q_a-a_{\rm L}\Pi_{\rm L}\succeq_{\vartheta_0}a_{\rm D}\Pi_{F_0}$, and
 under the root envelope and $a_{\rm D}\ge\delta a_{\rm L}$ of (b) the measure
 $\mathcal Q_a$ is reflection positive at every chamber cut. For
 $t_Y=\frac12$ the conditions
 follow from $H\le\frac32$, $M\ge2$, $s\ge q_\perp+M\rho_{\rm E}+2\beta^2$ and
 $kM\ge s^2$. For graded weights, with $h=\rho_{\rm E}/r$ and
 $f_{\rm al}=(q_\perp+\beta^2/t_Y)/r$, the conditions read
 $w_{\rm D}\ge d\,w_{\rm L}$, $w_{\rm S}\ge(H+t_Y)w_{\rm D}$,
 $w_{\rm R}\ge f_{\rm al}w_{\rm D}+hw_{\rm S}$ and $w_{\rm C}w_{\rm S}\ge w_{\rm R}^2$,
 together with $t_Y\le H^{-1}$.
\item \emph{Two coordinates (slice principle).} In the setting of
 \cref{def:two-coordinate-five-type}, suppose that for each coordinate $c$ and
 each field type $\mathfrak f'$ the row data of $c$ at $\mathfrak f'$ form a
 five-type datum, and that for each type $t'$ the five weights $a_t$ of the
 states $t$ with $t_{1-c}=t'$, as $t_c$ ranges over the five types, satisfy
 the hypotheses of (b), of (c) or of (f) for the row data of $c$ at
 $\mathfrak f(t')$, these five weights playing the roles of
 $a_{\rm L},a_{\rm D},a_{\rm S},a_{\rm R},a_{\rm C}$ there respectively, in
 each case together with the root envelope and the inequality
 $a_{\rm D}\ge\delta a_{\rm L}$ of (b) in these roles. Then
 $\mathcal Q$ is reflection positive at every chamber cut $\vartheta_{c,i}$ of
 both coordinates. More precisely, for fixed $c$, $t'$ and $\phi'\in P_{t'}$,
 the sum of the components $a_t\Pi_{t,\phi}$ with $t_{1-c}=t'$ and
 $\phi_{1-c}=\phi'$ is the image of the five-type measure of these row data
 under the emission (placement map) $e_{\phi'}\colon\mathfrak X^{\Z/6\Z}\to X^W$,
 $e_{\phi'}(\Psi)_{(u_0,u_1,\ell)}=\Psi_{u_c}(\phi'(u_{1-c}),\ell)$.
\item \emph{The reserved cut (fold pairing).} Let $\Lambda=\{0,1\}$, let $\iota$
 be the layer swap, and for a one-layer measure $\kappa$ let
 $\operatorname{Diag}_\iota(\kappa)$ be the law that carries the same one-layer
 configuration on both layers. Consider the components of one coordinate
 (objects carrying one label per layer) or of two coordinates as in (d), and
 the one-layer components $\Pi^{(1)}$ built in the same way from one-layer
 laws; let $\pi_{t,j}$ be the sum of the $j$-layer components of the type or
 state $t$. Suppose that every two-layer component $\Pi$ of $t$ satisfies
 $\Pi\succeq_\iota-k_tr\operatorname{Diag}_\iota(\Pi^{(1)})$, where
 $\Pi^{(1)}$ is a fixed one-layer component with the same patterns as
 $\Pi$ (one such component being fixed for each two-layer component), and
 let $w_t,v_t\ge0$ with
 $v_t\ge k_tw_t$ for every $t$. Then
 \[
  \widetilde\mu=\sum_tw_tr^{\deg t}\pi_{t,2}
   +r\sum_tv_tr^{\deg t}\operatorname{Diag}_\iota(\pi_{t,1})\succeq_\iota0 .
 \]
 $\operatorname{Diag}_\iota$ preserves the order $\succeq_\vartheta$ at every
 chamber cut $\vartheta$; in particular it maps a one-layer measure that is
 reflection positive at a chamber cut to a measure reflection positive there.
 Hence if the two-layer data with the weights $w_tr^{\deg t}$ and the one-layer
 data with the weights $v_tr^{\deg t}$ satisfy the hypotheses of (b), (c), (d)
 or (f), root envelopes included (for (b) in graded form: $v_{\rm D}\ge d_1v_{\rm L}$,
 $v_{\rm S}\ge v_{\rm D}$, $v_{\rm R}\ge f_1v_{\rm D}+h_1v_{\rm S}$,
 $v_{\rm C}v_{\rm S}\ge v_{\rm R}^2$ with the one-layer inputs; all these conditions
 are homogeneous in $v$), then $\widetilde\mu$ is
 reflection positive at every chamber cut of the active coordinates and at
 $\iota$. Finally, if a field
 law $\lambda$ and its one-layer law $\lambda^{(1)}$ satisfy
 $\lambda\succeq_\iota-kr\operatorname{Diag}_\iota(\lambda^{(1)})$, then so do
 their images under every emission $e_\phi$ (in one coordinate, the map
 $\omega\mapsto\omega\circ\phi$); so components that use only some
 objects of a field law, the merges and the constants, inherit its bound. A
 component may also be used with a smaller proved coefficient than $k_t$
 in its own order at $\iota$, and $\Pi^{(1)}$ need
 not be the one-layer marginal of $\Pi$: only the order
 $\Pi\succeq_\iota-k_tr\operatorname{Diag}_\iota(\Pi^{(1)})$ is used.
\item \emph{Joint Schur criterion.} Suppose that for every $b\in\supp\nu$ the
 matching kernel has exact marginals, as in (c), and the fold kernel has the
 \emph{joint data}
 \[
  K_b=T_{\rm F}^b=m_b\bigl(\varphi_0\otimes\varphi_2+C_b+R_b\bigr),\qquad
  m_b\in[H^{-1},H],
 \]
 with $\E\varphi_0=\E\varphi_2=1$ (both $\varphi_0$ and $\varphi_2$
 depend on $b$, the index being suppressed) and
 $\|\varphi_o-\one\|\le\upsilon$ ($o=0,2$),
 where $\varphi_0\otimes\varphi_2$ is the kernel $\varphi_0(x)\varphi_2(y)$ and the
 kernels $C_b$, $R_b$ annihilate the constants on both sides, with
 $\|R_b\|\le\eps_{\rm T}$; and that, with
 $E_b=\frac12(T_{\rm E}^b+(T_{\rm E}^b)^*)-\Jproj$,
 \[
  E_b\succeq-\eps_{\rm E}I,\qquad
  \begin{pmatrix}E_b&C_b\\ C_b^*&E_b\end{pmatrix}\succeq-\kappa_{\rm J}I
  \quad\text{on }L^2(\nu^b)\oplus L^2(\nu^b).
 \]
 Let $a_{\rm D}$, $M$, $s$ and $k$ be as in (c), and let $t_Y>0$ satisfy
 \[
  t_Y\le H^{-1},\qquad M\ge H+t_Y,\qquad
  s\ge s_{\rm J},\qquad kM\ge s^2,
 \]
 where
 \[
  s_{\rm J}=M\eps_{\rm E}+H(\kappa_{\rm J}-\eps_{\rm E})_++H\eps_{\rm T}
  +\Bigl(H+\frac{H^2}{t_Y}\Bigr)\upsilon^2 .
 \]
 Then the conclusions of (c) hold. For graded weights, with
 $h=\eps_{\rm E}/r$, $k_{\rm J}=\kappa_{\rm J}/r$ and
 \[
  f_{\rm J}=\frac{H\{(k_{\rm J}-h)_+r+\eps_{\rm T}\}+(H+H^2/t_Y)\upsilon^2}{r},
 \]
 the conditions read $w_{\rm D}\ge d\,w_{\rm L}$,
 $w_{\rm S}\ge(H+t_Y)w_{\rm D}$, $w_{\rm R}\ge f_{\rm J}w_{\rm D}+hw_{\rm S}$ and
 $w_{\rm C}w_{\rm S}\ge w_{\rm R}^2$, together with $t_Y\le H^{-1}$.
\end{enumerate}
\end{lemma}

\begin{proof}
(a) By \cref{lem:record:fold-identification}(e), precomposition with
$\alpha\in\mathcal G_6$ permutes the fold patterns, maps $E_s$ to
$\psi\circ E_s$, permutes the merges of each partition and fixes the constant
maps. The law of $\Psi\circ\alpha$ under $\Pi_\phi(\lambda)$ is
$\Pi_{\phi\circ\alpha}(\lambda)$, and $\Pi_{\psi\circ E_s}(\pi_{\rm E})=
\Pi_{E_s}(\pi_{\rm E})$ because $\pi_{\rm E}$ is exchangeable. Hence every sum
$\pi_t$ with $t\ne{\rm L}$ is invariant, and $\Pi_{\rm L}$ is invariant by (v).

The cut forms are read off \cref{tab:corner}. The halves of $\Pi_{F_0}$
coincide. For $F_1$ the halves are $(\omega_1,\omega_0,\omega_0)$ and
$(\omega_2,\omega_2,\omega_1)$, so its form is
$\E_{\pi_{\rm F}}[U_{\omega_1}(\omega_0)V_{\omega_1}(\omega_2)]=\E_b\ip{U_b}{T_{\rm F}^bV_b}$
by (iv); $F_2$ exchanges the two halves and gives the same value. For $E_1$
the form is $\E_{\pi_{\rm E}}[V_{\omega_1}(\omega_0)V_{\omega_1}(\omega_2)]=
\E_b\ip{V_b}{T_{\rm E}^bV_b}$ by (iii). For $E_0$ the object on both halves is
$\omega_0$, and the form is $\E_{\pi_{\rm E}}[U_{\omega_0}(\omega_1)U_{\omega_0}(\omega_2)]
=\E_b\ip{U_b}{T_{\rm E}^bU_b}$, because $(\omega_1,\omega_0,\omega_2)$ again has the
law $\pi_{\rm E}$. For the merges $(\omega_0,\omega_1)\sim\Gamma$: $G_e$ gives
$\E_\Gamma[U_{\omega_1}(\omega_0)^2]=\E_b\|U_b\|^2$ by (ii), and $G_{01}$, $G_{02}$
give $\E_\Gamma[z_{\omega_0}U_{\omega_0}(\omega_1)]=\E_b[z_bu_b]$ by (ii) and the
symmetry of $\Gamma$; $G_o$, $G_{10}$ and $G_{12}$ are the same with $V$ in
place of $U$. The constants give $\E_\nu z_b^2$. Finally
$\ip\Phi{T\Psi}=\ip\Phi{\Jproj\Psi}+\ip\Phi{(T-\Jproj)\Psi}$ and
$\ip\Phi{\Jproj\Psi}=\E\Phi\,\E\Psi$.

(b) Fix $b\in\supp\nu$ and omit it from the notation. Let $\Phi_b$ be the sum of
the terms of the integrand of $Q_{\mathcal Q_a-a_{\rm L}\Pi_{\rm L}-a_{\rm D}\Pi_{F_0}}(g)$
given by (a):
\[
 \Phi_b=2a_{\rm D}\ip U{T_{\rm F}V}+a_{\rm S}\bigl(\ip U{T_{\rm E}U}+\ip V{T_{\rm E}V}\bigr)
 +a_{\rm R}\bigl(\|U\|^2+\|V\|^2+2z(u+v)\bigr)+2a_{\rm C}z^2 .
\]
By (a) and $2\|U\|\|V\|\le\|U\|^2+\|V\|^2$,
\[
 \Phi_b\ge(a_{\rm R}-e_{\rm F}a_{\rm D}-e_{\rm E}a_{\rm S})\bigl(\|U\|^2+\|V\|^2\bigr)
 +a_{\rm S}(u^2+v^2)+2a_{\rm D}uv+2a_{\rm R}z(u+v)+2a_{\rm C}z^2 .
\]
The first term is nonnegative. With $w=u+v$,
$a_{\rm S}(u^2+v^2)+2a_{\rm D}uv=\frac12(a_{\rm S}+a_{\rm D})w^2
+\frac12(a_{\rm S}-a_{\rm D})(u-v)^2\ge\frac12a_{\rm S}w^2$, because
$a_{\rm S}\ge a_{\rm D}\ge0$. The binary form
$\frac12a_{\rm S}w^2+2a_{\rm R}zw+2a_{\rm C}z^2$ has nonnegative diagonal
coefficients and determinant $a_{\rm S}a_{\rm C}-a_{\rm R}^2\ge0$, so
$\Phi_b\ge0$. The cut form at $\vartheta_0$ of
$\mathcal Q_a-a_{\rm L}\Pi_{\rm L}-a_{\rm D}\Pi_{F_0}$ at $g$ is $\E_b\Phi_b\ge0$,
and its cut matrix is symmetric, because $\mathcal Q_a$ and $\Pi_{\rm L}$ are
$\vartheta_0$-invariant by (a) and $\Pi_{F_0}$ is fixed samplewise by
$\vartheta_0$ (\cref{lem:record:fold-identification}(b)). This is the first
claim. Under the root envelope and $a_{\rm D}\ge\delta a_{\rm L}$ (for
$a_{\rm L}=0$ the term $a_{\rm L}\Pi_{\rm L}$ vanishes, and the chain holds with
$\delta=0$),
\[
 \mathcal Q_a\succeq_{\vartheta_0}a_{\rm L}\Pi_{\rm L}+a_{\rm D}\Pi_{F_0}
 \succeq_{\vartheta_0}(a_{\rm D}-\delta a_{\rm L})\Pi_{F_0}\succeq_{\vartheta_0}0,
\]
the last step by \cref{lem:commoncut}(i). Since $\mathcal Q_a$ is
$\mathcal G_6$-invariant, chamber transport (\cref{lem:slot-array}(i))
gives every cut $\vartheta_i$. For the graded weights, divide
$a_{\rm D}\ge\delta a_{\rm L}$ and the first displayed condition by $r$ and the
second by $r^2$; in the third both sides carry $r^4$, since $3+1=2+2$.

(c) Fix $b$ and omit it. Write $U=u+U^\circ$, $V=v+V^\circ$ with
$U^\circ,V^\circ\perp\one$, $\sigma_U=\|U^\circ\|$, $\sigma_V=\|V^\circ\|$, and
$K=K_b$; only $m_b$ keeps its index. Then
$\ip U{KV}=m_buv+u\ip{(I-\Jproj)K^*\one}{V^\circ}+v\ip{U^\circ}{(I-\Jproj)K\one}
+\ip{U^\circ}{(I-\Jproj)K(I-\Jproj)V^\circ}$, so the fold data give
\[
 2\ip U{KV}\ge2m_buv-2\beta\bigl(|u|\sigma_V+|v|\sigma_U\bigr)
 -q_\perp(\sigma_U^2+\sigma_V^2).
\]
Since $T_{\rm E}$ has exact marginals,
$\ip U{T_{\rm E}U}=u^2+\ip{U^\circ}{(T_{\rm E}-\Jproj)U^\circ}\ge u^2-\rho_{\rm E}\sigma_U^2$,
and likewise for $V$; and $\|U\|^2=u^2+\sigma_U^2$. Hence $\Phi_b/a_{\rm D}$, with
$\Phi_b$ as in (b), is at least
\[
 (s-q_\perp-M\rho_{\rm E})(\sigma_U^2+\sigma_V^2)+(M+s)(u^2+v^2)+2m_buv
 -2\beta\bigl(|u|\sigma_V+|v|\sigma_U\bigr)+2sz(u+v)+2kz^2 .
\]
By $2\beta|u|\sigma_V\le t_Yu^2+\beta^2\sigma_V^2/t_Y$ and the same
for $(v,\sigma_U)$, this is at least
\[
 \Bigl(s-q_\perp-M\rho_{\rm E}-\frac{\beta^2}{t_Y}\Bigr)(\sigma_U^2+\sigma_V^2)
 +A(u^2+v^2)+2m_buv+2sz(u+v)+2kz^2,\qquad A=M+s-t_Y.
\]
The fluctuation term is nonnegative. With $w_\pm=(u\pm v)/\sqrt2$ the rest is
$(A-m_b)w_-^2+(A+m_b)w_+^2+2\sqrt2\,szw_++2kz^2$. Here
$A-m_b\ge M-t_Y-H+s\ge s\ge0$, and $A+m_b\ge M+s-t_Y+H^{-1}\ge M$, so
the binary form in $(w_+,z)$ has nonnegative diagonal coefficients and
determinant $2k(A+m_b)-2s^2\ge2(kM-s^2)\ge0$. Thus $\Phi_b\ge0$, and the rest of
the proof is that of (b). For $t_Y=\frac12$: $M\ge2\ge H+\frac12$ and
$\frac12<\frac23\le H^{-1}$. The graded form follows as in (b), since
$s=w_{\rm R}r/w_{\rm D}$, $M=w_{\rm S}/w_{\rm D}$ and $k=w_{\rm C}r^2/w_{\rm D}$.

(d) Fix $c=0$; the coordinate $1$ is the same with the coordinates exchanged.
Fix a type $t'$ of the coordinate $1$ and $\phi_1\in P_{t'}$, and let
$\mathcal Q^{\rm row}$ be the five-type measure of the row data of $0$ at
$\mathfrak f(t')$ with the weights $a_{(t_0,t')}$. For a type $t_0$ and
$\phi_0\in P_{t_0}$, the configuration $\Psi_{u_0}=y_{\phi_0(u_0),\cdot,\cdot}$ of
rows of $y\sim\lambda_{\mathfrak f(t_0)\mathfrak f(t')}$ has the law of the
component $\phi_0$ of the row data: the rows of $\lambda_{{\rm L}\mathfrak f(t')}$
are $\pi_{\rm L}$, the three rows of $\lambda_{{\rm F}\mathfrak f(t')}$ are
$\pi_{\rm F}$, those of $\lambda_{{\rm E}\mathfrak f(t')}$ are $\pi_{\rm E}$, and
a merge or a constant reads only the rows $0,1$ or the row $0$, whose laws
are $\Gamma$ and $\nu$. Moreover
$e_{\phi_1}(\Psi)_{(u_0,u_1,\ell)}=y_{\phi_0(u_0),\phi_1(u_1),\ell}$. So the sum
of $a_t\Pi_{t,\phi}$ over the states $t=(t_0,t')$ and the patterns
$(\phi_0,\phi_1)$ is $(e_{\phi_1})_*\mathcal Q^{\rm row}$, which is the slice
decomposition. Let $\vartheta=\vartheta_{0,i}$, with plus half
$W_+=H_i\times\Z/6\Z\times\Lambda$. The emission is half-local for the cuts
$\vartheta_i$ of $\Z/6\Z$ and $\vartheta$ of $W$ in the sense of
\cref{lem:commoncut}(v): $e_{\phi_1}(\Psi)|_{W_+}=\varphi(\Psi|_{H_i})$ and the
reflected half of $e_{\phi_1}(\Psi)$ is $\varphi$ of the reflected half of
$\Psi$, with
$\varphi(x)=(x_{u_0}(\phi_1(u_1),\ell))_{(u_0,u_1,\ell)\in W_+}$ for
$x\in\mathfrak X^{H_i}$; so the image of a measure reflection positive at
$\vartheta_i$ is reflection positive at $\vartheta$. By (b), (c) or (f),
$\mathcal Q^{\rm row}$ is reflection positive at every $\vartheta_i$, and so is
each image. Summing over $t'$ and $\phi_1$ proves the claim.

(e) A component $\Pi$ of $t$, with the weight $w_tr^{\deg t}$, together with
its one-layer diagonal, with the weight $v_tr^{\deg t+1}$, satisfies
\[
 w_tr^{\deg t}\Pi+v_tr^{\deg t+1}\operatorname{Diag}_\iota(\Pi^{(1)})
 \succeq_\iota(v_t-k_tw_t)\,r^{\deg t+1}\operatorname{Diag}_\iota(\Pi^{(1)})
 \succeq_\iota0,
\]
by the hypothesis and because a diagonal copy is fixed samplewise by $\iota$
(\cref{lem:commoncut}(i)); summing over the components gives
$\widetilde\mu\succeq_\iota0$. Let $\vartheta$ be a chamber cut, acting on the
chamber coordinates only. The diagonal embedding $z'\mapsto z$,
$z_{(u,\ell)}=z'_u$, is half-local for $\vartheta$ on the one-layer and on the
two-layer slots, so by \cref{lem:commoncut}(v) it preserves the order
$\succeq_\vartheta$; this is the statement about $\operatorname{Diag}_\iota$, and
the two sums of $\widetilde\mu$
are then reflection positive at the chamber cuts by (b), (c), (d) or (f). Finally
$e_\phi$ is half-local for the layer swaps of the field labels and of the
slots, so by \cref{lem:commoncut}(v) the image of
$\lambda+kr\operatorname{Diag}_\iota(\lambda^{(1)})\succeq_\iota0$ is
$(e_\phi)_*\lambda+kr\operatorname{Diag}_\iota((e_\phi)_*\lambda^{(1)})\succeq_\iota0$.
The argument uses each component only through its own order at $\iota$ and
the one-layer sums only through their own completion inequalities; this is
the last sentence of (e).

(f) Fix $b$ and omit it; write $U=u+U^\circ$, $V=v+V^\circ$, $\sigma_U$, $\sigma_V$
as in (c), and $\alpha=\ip{U^\circ}{\varphi_0-\one}$,
$\alpha'=\ip{\varphi_2-\one}{V^\circ}$, so that $|\alpha|\le\upsilon\sigma_U$ and
$|\alpha'|\le\upsilon\sigma_V$. Since $\E\varphi_0=\E\varphi_2=1$ and $C_b$, $R_b$
annihilate the constants on both sides,
\[
 \ip U{K_bV}=m_b(u+\alpha)(v+\alpha')+m_b\ip{U^\circ}{(C_b+R_b)V^\circ}.
\]
By $0\le m_b\le H$ and Young's inequality with the allowance $t_Y$,
\begin{gather*}
 2m_b(u\alpha'+v\alpha)\ge-t_Y(u^2+v^2)
 -\frac{H^2\upsilon^2}{t_Y}(\sigma_U^2+\sigma_V^2),\\
 2m_b\bigl(\alpha\alpha'+\ip{U^\circ}{R_bV^\circ}\bigr)\ge
 -H(\upsilon^2+\eps_{\rm T})(\sigma_U^2+\sigma_V^2).
\end{gather*}
The matching kernel has exact marginals, so
$\ip U{T_{\rm E}U}=u^2+\ip{U^\circ}{E_bU^\circ}$, and likewise for $V$. The
remaining terms $a_{\rm S}(\ip{U^\circ}{E_bU^\circ}+\ip{V^\circ}{E_bV^\circ})
+2m_ba_{\rm D}\ip{U^\circ}{C_bV^\circ}$ are the form of the block operator
$\bigl(\begin{smallmatrix}a_{\rm S}E_b&m_ba_{\rm D}C_b\\
m_ba_{\rm D}C_b^*&a_{\rm S}E_b\end{smallmatrix}\bigr)$ at $(U^\circ,V^\circ)$. Apply \cref{cor:joint-row-split} with
$h=\eps_{\rm E}$, $\mathsf k_\Blk=\max(\kappa_{\rm J},\eps_{\rm E})$ and
$m=m_b$. Its hypotheses hold because $a_{\rm S}=Ma_{\rm D}\ge Ha_{\rm D}$
and $m_b\le H$; the splitting uses only these block inequalities, not the
prime-product origin of the operators. Thus the remaining terms are at least
\[
 -\bigl(a_{\rm S}\eps_{\rm E}+Ha_{\rm D}(\kappa_{\rm J}-\eps_{\rm E})_+\bigr)
 (\sigma_U^2+\sigma_V^2).
\]
Collecting, $\Phi_b/a_{\rm D}$, with $\Phi_b$ as in (b), is at least
\[
 (s-s_{\rm J})(\sigma_U^2+\sigma_V^2)+A(u^2+v^2)+2m_buv+2sz(u+v)+2kz^2,
 \qquad A=M+s-t_Y,
\]
which is the form of the proof of (c) with the fluctuation coefficient
$s-s_{\rm J}\ge0$. Its mean part is nonnegative exactly as there, because
$M\ge H+t_Y$, $t_Y\le H^{-1}$ and $kM\ge s^2$; so $\Phi_b\ge0$, and the
rest of the proof is that of (b). In graded form $\eps_{\rm E}=hr$ and
$\kappa_{\rm J}=k_{\rm J}r$, so $s_{\rm J}=(w_{\rm S}/w_{\rm D})hr+f_{\rm J}r$ and
$s\ge s_{\rm J}$ reads $w_{\rm R}\ge f_{\rm J}w_{\rm D}+hw_{\rm S}$; the other conditions
read as in (c).
\end{proof}

\begin{corollary}[Joint rows of a block]\label{cor:joint-completion-blocks}
Let $\Blk$ be a finite set of primes with least prime $P$, $r=P^{-1/2}$ and
$r_p=p^{-1/2}$. Fix a value $b$ of the middle object of a five-type datum whose
conditional kernels at $b$ are products over the primes of $\Blk$:
$\nu^b=\bigotimes_p\nu_p$, $T^b_{\rm E}=\bigotimes_pT_{{\rm E},p}$ with
$T_{{\rm E},p}$ a self-adjoint coupling of $\nu_p$ with itself, and
$T^b_{\rm F}=m_bT_{{\rm F},\Blk}$, $T_{{\rm F},\Blk}=\bigotimes_pT_{{\rm F},p}$, with
$T_{{\rm F},p}$ of mass one and marginal densities $u_{{\rm L},p}$, $u_{{\rm R},p}$
relative to $\nu_p$. Let $C_{{\rm F},p}=T_{{\rm F},p}-u_{{\rm L},p}\otimes u_{{\rm R},p}$
and $C_{{\rm E},p}=T_{{\rm E},p}-\Jproj_p$, as in \cref{lem:joint-rows}, and suppose
that for every $p\in\Blk$
\[
 \|C_{{\rm F},p}\|\le c_{\rm F}r_p,\qquad \|C_{{\rm E},p}\|\le c_{\rm E}r_p,\qquad
 \begin{pmatrix}C_{{\rm E},p}&C_{{\rm F},p}\\ C_{{\rm F},p}^*&C_{{\rm E},p}\end{pmatrix}
 \succeq-k_{\rm J}r_pI,
\]
with $c_{\rm F}r,c_{\rm E}r<1$ and $(c_{\rm F}^2+c_{\rm E}^2)r\le k_{\rm J}$. Let
$m_b\in[H^{-1},H]$, let the block marginal densities
$U_{\rm L}=\bigotimes_pu_{{\rm L},p}$ and $U_{\rm R}=\bigotimes_pu_{{\rm R},p}$ satisfy
$\|U_{\rm L}-\one\|,\|U_{\rm R}-\one\|\le\upsilon$, let
$E_\Blk=T^b_{\rm E}-\Jproj\succeq-hrI$, and put
\[
 E_{\rm T}=(2\upsilon+\upsilon^2)\Bigl(\prod_{p\in\Blk}(1+c_{\rm F}r_p)-1\Bigr),
\]
which is at most $(2\upsilon+\upsilon^2)(e^{nc_{\rm F}r}-1)$ if $\Blk$ has at most
$n$ primes. Then the joint data of \cref{lem:five-type}(f) hold at $b$ with
$\varphi_0=U_{\rm L}$, $\varphi_2=U_{\rm R}$, $\eps_{\rm E}=hr$,
$\kappa_{\rm J}=k_{\rm J}r$ and $\eps_{\rm T}=E_{\rm T}$. Consequently, if this holds
at every $b\in\supp\nu$ and $0<t_Y\le H^{-1}$, then for row parameters
\[
 M\ge H+t_Y,\qquad
 f\ge\frac{H\{(k_{\rm J}-h)_+r+E_{\rm T}\}+(H+H^2/t_Y)\upsilon^2}{r},
\]
and $h$, the completion inequalities $w_{\rm D}\ge d\,w_{\rm L}$,
$w_{\rm S}\ge Mw_{\rm D}$, $w_{\rm R}\ge fw_{\rm D}+hw_{\rm S}$,
$w_{\rm C}w_{\rm S}\ge w_{\rm R}^2$ imply the graded conditions of
\cref{lem:five-type}(f). This holds for the two-layer and for the one-layer rows
alike, so the corollary serves the ordinary completion, with the two-layer rows,
and the reserved one, with the one-layer rows (\cref{lem:five-type}(e)).
\end{corollary}

\begin{proof}
The kernels $C_{{\rm F},p}$ and $C_{{\rm E},p}$ annihilate the constants on both
sides. Hence $C_\Blk=\bigotimes_p(\Jproj_p+C_{{\rm F},p})-\Jproj_\Blk$ and
$E_\Blk=\bigotimes_p(\Jproj_p+C_{{\rm E},p})-\Jproj_\Blk$ annihilate them too, and
their joint block is bounded by \cref{lem:joint-rows}(b). Indeed on its
singleton sectors the joint block is $\succeq-k_{\rm J}rI$ (using
$r_p\le r$), and on sectors with at least two primes it is
$\succeq-(c_{\rm F}^2+c_{\rm E}^2)r^2I\succeq-k_{\rm J}rI$.
Hence the joint block of $(E_\Blk,C_\Blk)$ is $\succeq-k_{\rm J}rI$.
Write
$T_{{\rm F},\Blk}=U_{\rm L}\otimes U_{\rm R}+C_\Blk+R_\Blk$. Since
$T_{{\rm F},\Blk}\one=U_{\rm L}$, $T_{{\rm F},\Blk}^*\one=U_{\rm R}$ and $C_\Blk$
annihilates the constants, so does $R_\Blk$; and
$\|R_\Blk\|\le(2\upsilon+\upsilon^2)(\prod_p(1+\|C_{{\rm F},p}\|)-1)\le E_{\rm T}$ by
\cref{lem:joint-rows}(c). With $C_b=C_\Blk$ and $R_b=R_\Blk$ these are the joint
data of \cref{lem:five-type}(f); the bound by $e^{nc_{\rm F}r}$ is $1+x\le e^x$.
Finally $f\ge f_{\rm J}$ and $M\ge H+t_Y$, so the four completion
inequalities with $(d,M,f,h)$ imply those of \cref{lem:five-type}(f). Nothing in
the argument depends on the number of layers.
\end{proof}

\begin{remark}[Fold data from marginal densities]\label{rem:five-type-fold-data}
Suppose that, for every $b\in\supp\nu$, the middle density $m_b$ --- the
density at $b$ of the $\omega_1$-marginal of $\pi_{\rm F}$ relative to
$\nu$, equal to $\ip\one{K_b\one}$ for $K_b=T_{\rm F}^b$ by the
factorisation in \cref{def:five-type-datum}(iv) --- lies in $[H^{-1},H]$, and the conditional law of $(\omega_0,\omega_2)$ given
$\omega_1=b$ under $\pi_{\rm F}$ has marginal densities $\varphi_0,\varphi_2$
relative to $\nu^b$, with $\varphi_o\le R_{\rm m}$ and
$\|\varphi_o-\one\|\le\upsilon$, and maximal correlation at most $\rho_{\rm F}$
relative to its own marginals. Then the fold data of (c) hold with
$\beta=H\upsilon$ and $q_\perp=H(R_{\rm m}\rho_{\rm F}+\upsilon^2)$, and
\[
 f_{\rm al}=\frac{HR_{\rm m}\rho_{\rm F}+(H+H^2/t_Y)\upsilon^2}{r} .
\]
Indeed $K_b=m_b(\varphi_0\otimes\varphi_2+c_b)$, where the kernel $c_b$ represents
the covariance, $\ip h{c_bh'}=\operatorname{Cov}(h(\omega_0),h'(\omega_2)\mid\omega_1=b)$;
so $c_b\one=c_b^*\one=0$ and $\|c_b\|\le R_{\rm m}\rho_{\rm F}$, because
$\|h\|^2_{\varphi_o\nu^b}\le R_{\rm m}\|h\|^2$. Hence $\ip\one{K_b\one}=m_b$,
$(I-\Jproj)K_b\one=m_b(\varphi_0-\one)$, and the perpendicular block is
$m_b\bigl((\varphi_0-\one)\otimes(\varphi_2-\one)+(I-\Jproj)c_b(I-\Jproj)\bigr)$. With
$t_Y=1/50$ this is the Schur row coefficient
$f_{\rm Schur}=(HR_{\rm m}\rho_{\rm F}+(H+50H^2)\upsilon^2)/r$ of the tail datum of the
upper range (\ref{it:record:upper-F5}). The joint form (f) splits the
covariance kernel as $c_b=C_b+R_b$ and bounds $C_b$ together with the matching
row, in one Gram matrix, instead of adding the two norms: the term
$HR_{\rm m}\rho_{\rm F}$ becomes $H\{(k_{\rm J}-h)_+r+\eps_{\rm T}\}$, and on a block
with the data of \cref{cor:joint-completion-blocks}
\[
 f_{\rm J}=\frac{H\{(k_{\rm J}-h)_+r+E_{\rm T}\}+(H+H^2/t_Y)\upsilon^2}{r},
\]
with a rational $t_Y$ chosen for each row; this is the row coefficient of the
bands of the upper range (\ref{it:record:upper-F5}). Both criteria keep a
change of marginals and a centred correlation apart: in the constant direction the
matching contribution dominates the fold term because $M\ge H+t_Y$ while
$m_b\le H$, and Young's inequality with the allowance $t_Y$ absorbs the
constant--perpendicular coupling of the fold, so a linear marginal error
enters only through its square. The allowance enters only through
$M\ge H+t_Y$, $t_Y\le H^{-1}$ and the term $H^2\upsilon^2/t_Y$, so
it may be chosen freely, row by row.
\end{remark}

\begin{remark}[The model computation and its ancestors]\label{rem:five-type-model}
The proof of (b), (c) and (f) is one computation: the cut forms of the twelve
components other than the root and the aligned fold are explicit in the
conditional means $u_b,v_b$, the fluctuations and the constant value $z_b$, and
the four inequalities are the conditions for the resulting form in $(u,v,z)$ to
be positive after the rotation $w_\pm=(u\pm v)/\sqrt2$. Only the inequalities
matter: any nonnegative weights satisfying them give a positive law, so the
certificates may choose rational weights by a sequential recipe or jointly, and
then check the inequalities exactly. Without a root ($a_{\rm L}=0$) and with
$a_{\rm D}=1$, $a_{\rm S}=1+a_{\rm R}$ and $a_{\rm C}=a_{\rm R}$, (b) is the array
completion of the middle range of the $216$-slot design: its condition
$a_{\rm R}(1-e_{\rm E})\ge e_{\rm F}+e_{\rm E}$ is the second inequality of (b),
and the other two hold for these weights. Its later row completions, with a
separate bound for the middle density and with split marginal and fluctuation
bounds, are described in \cref{sec:history}.
\end{remark}

\begin{remark}[How the completions are certified]\label{rem:five-type-certificates}
Every certified completion inequality is an instance of (b)--(f) with the
stored inputs. A middle reference (\cref{lem:front12:reference}) uses
(b) with the full-norm row coefficients $f_2,h_2$, (b) for the one-layer
diagonals with $f_1,h_1$, and (e): $4+4+5=13$ inequalities per piece, $429$ for
the $33$ pieces (\cref{cert:record:frame-data}(c)); the type-S inequality of (e) (the matching type, of field
type ${\rm E}$) is checked in the stronger form
$v_{\rm S}\ge2k_{\rm E}w_{\rm S}$, where $k_{\rm E}=k_{\rm S}$ is the
coefficient $k_t$ of (e) for $t={\rm S}$, and the
protection inequality $v_{\rm D}\ge2k_{\rm D}w_{\rm D}$ of
\cref{lem:front12:reference}(iv), which implies the type-D inequality of (e), is
checked with them. An upper band uses (d) in both coordinates and for both
layer counts, with (e): $40$ slice inequalities for $w$, $40$ for $v$ and the $25$ domination
inequalities $v_t\ge k_tw_t$ of (e), one for each of the $25$ states,
$105$ per band and $32550$ for the $310$ bands
(\cref{cert:upper-completion}). Its slices use (f), through
\cref{cor:joint-completion-blocks}, with $h=\eps_{\rm E}/r$ and a rational mean
allowance $t_Y$ for each row, with $t_YH<1$ and $M\ge H+t_Y$. The
tail datum of the upper range (\cref{def:record:upper-bands}(e)) has its own
$105$ inequalities, and its slices use (c) with $t_Y=1/50$ and
$h=\rho_{\rm E}/r$: in all six rows of the tail datum (three field types, two
layer counts) the coefficient $f_{\rm Schur}$ of \cref{rem:five-type-fold-data} is
smaller than the full-norm one and $H<103/100$, and the floor
$w_{\rm S}\ge\frac{21}{20}w_{\rm D}$ of its slices gives $M\ge H+t_Y$ and
$t_Y\le H^{-1}$ (\ref{it:record:upper-F5}).
\end{remark}

\begin{remark}[Fold pairing]\label{rem:fold-pairing}
Part (e) pairs each component with the diagonal copy of a one-layer law, one
power of $r$ higher: the copy is fixed samplewise by the cut, hence positive
there, and it is positive at every other cut at which that law is. The
reserved cut is cheap because the half graph of a two-layer labelling at
$\iota$ is its one-layer graph, so the copy that cancels a component's negative
envelope is the one-layer version of the same component, and the one-layer
versions are made positive at the chamber cuts by their own completion
inequalities. The paired law need not be the half law of the component. On the
three late middle pieces the two-layer root is a positive reweighting of the
uniform root, and its layer cut is capped against the diagonal copy of the
uniform one-layer root, which the completion supplies
(\cref{sec:positive-roots}); the reserved merges and constants of the upper
range carry smaller caps, proved for the marginals they observe
(\cref{prop:observed-caps}). The same pairing, applied in every chamber
coordinate at once, made the high range of the $216$-slot design positive
(\cref{sec:history}).
\end{remark}

\subsection{Polarization: dominating two summands by two others}\label{subsec:polarization}

The first stage $\lambda$ must drop the fold records $F_1,F_2$,
which do not accept the queries, without losing positivity at the
chamber cut $\vartheta_0$ used for the first stage. At $\vartheta_0$
these two folds are nonaligned: their cut forms couple the two
corners $\mathcal U$ and $\mathcal V$ of the half space, while the
matching records $E_0,E_1$ act inside each corner separately
(\cref{tab:corner}). Polarization is the inequality that lets the
matchings replace the folds: the sum of the two matching forms
exceeds twice the fold form by the square $(u-v)^2$ of the
difference of the conditional means on the corners (a nonnegative
deficit term, favourable to us), less the explicit error term of
\eqref{eq:polarization} (the row errors), and those errors are in turn
dominated by the two merge records $G_e,G_o$, which are diagonal
on the corners. When the domination holds with both signs in
front of the fold terms, it also holds after taking tensor
products over any set of blocks (part~(d)), which the mean credit
of (CP7) uses.

\begin{lemma}[Polarization]\label{lem:tr:polarization}
Let $\nu^b$ be a probability measure on a finite set, and $T_{\rm F}$,
$T_{\rm E}$, $T_{\rm E'}$ kernels on $L^2(\nu^b)$ with
$\|T_{\rm F}-\Jproj\|\le e_{\rm F}$ and
$\|T_{\rm E}-\Jproj\|,\|T_{\rm E'}-\Jproj\|\le e_{\rm E}$. For real $U,V\in L^2(\nu^b)$
with means $u,v$,
\begin{equation}\label{eq:polarization}
 \ip U{T_{\rm E}U}+\ip V{T_{\rm E'}V}-2\ip U{T_{\rm F}V}\ge(u-v)^2
 -(e_{\rm F}+e_{\rm E})\bigl(\|U\|^2+\|V\|^2\bigr).
\end{equation}
Consequently, for the components of a five-type datum with
$\|T_{\rm F}^b-\Jproj\|\le e_{\rm F}$ and $\|T_{\rm E}^b-\Jproj\|\le e_{\rm E}$ for
every $b\in\supp\nu$, and every $\varepsilon\ge e_{\rm F}+e_{\rm E}$:
\begin{enumerate}[label=\textup{(\alph*)}]
\item $\Pi_{E_0}+\Pi_{E_1}+\varepsilon(\Pi_{G_e}+\Pi_{G_o})
 \succeq_{\vartheta_0}\Pi_{F_1}+\Pi_{F_2}$;
\item if $\nu^b(\{b\})=0$ for every $b\in\supp\nu$, so that the
 conditional law $\nu^b$ gives mass $0$ to the value $b$ itself and the
  middle object differs from each outer object almost surely under the
  fold law, then both orders
 $\Pi_{E_0}+\Pi_{E_1}+\varepsilon(\Pi_{G_e}+\Pi_{G_o})
 \succeq_{\vartheta_0}\pm(\Pi_{F_1}+\Pi_{F_2})$ hold. Equivalently, with $A=C_{\vartheta_0}(\Pi_{E_0}+\varepsilon\Pi_{G_e})$,
 $D=C_{\vartheta_0}(\Pi_{E_1}+\varepsilon\Pi_{G_o})$ and
 $K=C_{\vartheta_0}(\Pi_{F_1})$, viewed as matrices on
 $\R^{\mathfrak X^{H_0}}$ supported on the blocks
 $\mathcal U\times\mathcal U$, $\mathcal V\times\mathcal V$ and
 $\mathcal U\times\mathcal V$ respectively and zero elsewhere, both
 matrices
 $\bigl(\begin{smallmatrix}A&\pm K\\ \pm K^{\mathsf T}&D\end{smallmatrix}\bigr)$ on
 $\R^{\mathcal U}\oplus\R^{\mathcal V}$ are positive semidefinite; in
 particular $A,D\succeq0$;
\item let $X$ be a finite nonempty set of blocks, and for $\Blk\in X$ let a
 five-type datum on $\mathfrak X_{\Blk}$ satisfy the hypotheses of (b) with
 constants $\varepsilon_{\Blk}\ge e_{{\rm F},\Blk}+e_{{\rm E},\Blk}$. Then, as
 measures on $(\prod_{\Blk}\mathfrak X_{\Blk})^{\Z/6\Z}$ and for both signs,
 \[
  \bigotimes_{\Blk}\bigl(\Pi_{E_0,\Blk}+\varepsilon_{\Blk}\Pi_{G_e,\Blk}\bigr)
  +\bigotimes_{\Blk}\bigl(\Pi_{E_1,\Blk}+\varepsilon_{\Blk}\Pi_{G_o,\Blk}\bigr)
  \pm\Bigl(\bigotimes_{\Blk}\Pi_{F_1,\Blk}+\bigotimes_{\Blk}\Pi_{F_2,\Blk}\Bigr)
  \succeq_{\vartheta_0}0,
 \]
 and $\bigotimes_{\Blk}\bigl(\Pi_{E_0,\Blk}+\Pi_{E_1,\Blk}+\varepsilon_{\Blk}
 (\Pi_{G_e,\Blk}+\Pi_{G_o,\Blk})\bigr)\succeq_{\vartheta_0}
 \bigotimes_{\Blk}\Pi_{F_1,\Blk}+\bigotimes_{\Blk}\Pi_{F_2,\Blk}$;
\item for every number $k$ of tensor factors: if $\vartheta$ is a
 cut of a slot set $V$ and, for
 $j=1,\dots,k$, $\kappa_j$ and $\zeta_j$ are finite signed measures on
 $\mathfrak X_j^V$ with $\kappa_j\pm\zeta_j\succeq_\vartheta0$ for both signs, then
 \[
  \bigotimes_{j=1}^k\kappa_j-\bigotimes_{j=1}^k\zeta_j
  =2^{-k}\sum_{\sigma\in\{\pm1\}^k}(1-\sigma_1\cdots\sigma_k)
  \bigotimes_{j=1}^k(\kappa_j+\sigma_j\zeta_j),
 \]
 every product on the right is $\succeq_\vartheta0$, and
 $\bigotimes_j\kappa_j\pm\bigotimes_j\zeta_j\succeq_\vartheta0$ for both signs. In
 particular, under the hypotheses of (c), with
 $\mathcal Q^{\rm E}_\Blk=\Pi_{E_0,\Blk}+\Pi_{E_1,\Blk}
 +\varepsilon_\Blk(\Pi_{G_e,\Blk}+\Pi_{G_o,\Blk})$ and
 $\mathcal Q^{\rm F}_\Blk=\Pi_{F_1,\Blk}+\Pi_{F_2,\Blk}$,
 \[
  \bigotimes_{\Blk\in X}\mathcal Q^{\rm E}_\Blk\pm\bigotimes_{\Blk\in X}\mathcal Q^{\rm F}_\Blk
  \succeq_{\vartheta_0}0\qquad\text{for both signs.}
 \]
\end{enumerate}
\end{lemma}

\begin{proof}
By \cref{lem:five-type}(a),
$\ip U{T_{\rm E}U}\ge u^2-e_{\rm E}\|U\|^2$, $\ip V{T_{\rm E'}V}\ge v^2-e_{\rm E}\|V\|^2$
and $2\ip U{T_{\rm F}V}\le2uv+2e_{\rm F}\|U\|\|V\|\le2uv+e_{\rm F}(\|U\|^2+\|V\|^2)$,
which gives \eqref{eq:polarization}.

(a) By \cref{lem:five-type}(a), the cut form at $\vartheta_0$ of the difference of
the two sides at $g$ is
$\E_b\bigl[\ip{U_b}{T_{\rm E}^bU_b}+\ip{V_b}{T_{\rm E}^bV_b}+\varepsilon(\|U_b\|^2+\|V_b\|^2)
-2\ip{U_b}{T_{\rm F}^bV_b}\bigr]$, which is at least
$\E_b(u_b-v_b)^2\ge0$ by \eqref{eq:polarization}. The cut matrix is symmetric:
$\Pi_{E_s}$ and $\Pi_{G_e},\Pi_{G_o}$ are $\vartheta_0$-invariant, and
$\vartheta_0$ exchanges $\Pi_{F_1}$ and $\Pi_{F_2}$
(\cref{lem:record:fold-identification}(e)).

(b) If $\nu^b(\{b\})=0$, the fold law has
$\omega_1\ne\omega_0,\omega_2$ almost surely. The matching law has all three
objects distinct by exchangeability, and the merge law has its two objects
distinct. In particular, for $E_0$,
$\omega_0\ne\omega_2$ because $(\omega_0,\omega_2)$ has the law of
$(\omega_0,\omega_1)$ by exchangeability. Thus the
cut matrices of $\Pi_{E_0}$ and $\Pi_{G_e}$ live on $\mathcal U\times\mathcal U$,
those of $\Pi_{E_1}$ and $\Pi_{G_o}$ on $\mathcal V\times\mathcal V$, that of
$\Pi_{F_1}$ on $\mathcal U\times\mathcal V$ and that of $\Pi_{F_2}$, which is
$K^{\mathsf T}$ because $\Pi_{F_2}$ is the image of $\Pi_{F_1}$ under
$\vartheta_0$ (\cref{rem:cut-symmetries}(a)), on $\mathcal V\times\mathcal U$. Given $g$, let $\tilde g$ equal
$g$ on $\mathcal U$ and $-g$ on $\mathcal V$; the corners are disjoint
(\cref{lem:record:fold-identification}(f)). The forms of $\Pi_{E_s}$, $\Pi_{G_e}$,
$\Pi_{G_o}$ at $\tilde g$ equal those at $g$, and the fold forms change sign,
so (a) at $\tilde g$ is the order with the lower sign at $g$. The two orders are the two
block matrices evaluated at $(x|_{\mathcal U},x|_{\mathcal V})$, and a diagonal
block of a positive semidefinite matrix is positive semidefinite.

(c) Let $\mathcal H_{\Blk}=\R^{\mathfrak X_{\Blk}^{H_0}}$, and extend $A_{\Blk}$,
$D_{\Blk}$, $K_{\Blk}$ of (b) by zero to $\mathcal H_{\Blk}$. The block matrices
$\bigl(\begin{smallmatrix}A_{\Blk}&\pm K_{\Blk}\\ \pm K_{\Blk}^{\mathsf T}&D_{\Blk}
\end{smallmatrix}\bigr)$ remain positive semidefinite on
$\mathcal H_{\Blk}\oplus\mathcal H_{\Blk}$, because their form at $(x,y)$ is the
form of (b) at $(x|_{\mathcal U},y|_{\mathcal V})$. Choose the sign $-$ at one
block and $+$ at the others, or $+$ at every block. By
\cref{lem:record:block-gram},
$\bigl(\begin{smallmatrix}\bigotimes A_{\Blk}&\pm\bigotimes K_{\Blk}\\
\pm\bigotimes K_{\Blk}^{\mathsf T}&\bigotimes D_{\Blk}\end{smallmatrix}\bigr)\succeq0$
for both signs, and its form at $(x,x)$ is
$\ip x{(\bigotimes A_{\Blk}+\bigotimes D_{\Blk}\pm(\bigotimes K_{\Blk}+\bigotimes
K_{\Blk}^{\mathsf T}))x}$. By \eqref{eq:cc:kron} these four Kronecker products
are the cut matrices at $\vartheta_0$ of the four tensor products in the first
claim, which is therefore proved; the cut matrices are symmetric, each factor
being $\vartheta_0$-invariant or exchanged with its partner. For the second
claim expand the left side into the $2^{|X|}$ products
$\bigotimes_{\Blk}(\Pi_{E_{s_{\Blk}},\Blk}+\varepsilon_{\Blk}\Pi_{G_{s_{\Blk}s_{\Blk}},\Blk})$,
$s\in\{0,1\}^X$. The two constant choices of $s$ give the first claim with the
sign $-$; every other product has the cut matrix $\bigotimes_{\Blk}M_{\Blk}$
with $M_{\Blk}\in\{A_{\Blk},D_{\Blk}\}$, a Kronecker product of positive
semidefinite matrices.

(d) Expand the products on the right. The product with $\zeta_j$ at the
positions $j\in S$ and $\kappa_j$ elsewhere has the coefficient
$2^{-k}\sum_\sigma(1-\prod_j\sigma_j)\prod_{j\in S}\sigma_j=[S=\emptyset]-[S=\{1,\dots,k\}]$,
because $\sum_\sigma\prod_{j\in T}\sigma_j=2^k[T=\emptyset]$ for every
$T\subseteq\{1,\dots,k\}$ and $\prod_j\sigma_j\prod_{j\in S}\sigma_j=\prod_{j\notin S}\sigma_j$;
this is the identity. Each factor $\kappa_j+\sigma_j\zeta_j$ is
$\succeq_\vartheta0$, so each product is $\succeq_\vartheta0$ by \eqref{eq:cc:kron},
and the coefficients $1-\sigma_1\cdots\sigma_k\in\{0,2\}$ are nonnegative; hence
$\bigotimes_j\kappa_j-\bigotimes_j\zeta_j\succeq_\vartheta0$. Replacing $\zeta_1$ by
$-\zeta_1$, which satisfies the same hypothesis, gives the other sign. For cut
matrices this is \cref{lem:record:block-gram}(d); the identity moreover writes
the difference as a nonnegative combination of tensor products
of the $\succeq_\vartheta0$ signed measures
$\kappa_j+\sigma_j\zeta_j$, so that any of them may be discarded. For the last claim, (b) gives
$\mathcal Q^{\rm E}_\Blk\pm\mathcal Q^{\rm F}_\Blk\succeq_{\vartheta_0}0$ for every
$\Blk\in X$.
\end{proof}

\begin{remark}[Where polarization is used]\label{rem:polarization-uses}
(1) In the middle frame (\cref{def:record:frame}) (a) is applied to the
one-layer datum with $\varepsilon=(f_1+h_1)r$; multiplied by the record weight
$v_{\rm D}r^2$ and copied to the reserved diagonal, which preserves the order at
$\vartheta_0$ (\cref{lem:five-type}(e)), it is the first-stage order
$E_{\rm add}+G\succeq_{\vartheta_0}F_{\rm fold}$ of (CP2)(a). (2) In the variance
credits (a) is applied to the two-layer datum with
$\varepsilon_{\Blk}=(f_2+h_2)r_{\Blk}$ (\cref{lem:record:paired-e-bridge}(a)); this
datum also satisfies the hypothesis $\nu^b(\{b\})=0$ of (b), and (c)
gives the aligned bridge of several blocks (\cref{lem:record:paired-e-joint}(a)).
(3) The one-layer datum of (1) also satisfies the hypothesis of (b), whose two
orders, copied to the reserved diagonal as in (1), give
$E_{\rm add}+G\pm F_{\rm fold}\succeq_{\vartheta_0}0$ on every block; (d) carries
both orders to the products over any set of blocks, in every degree. This is
the fold coupon of the mean credit (\cref{prop:record:mean-credit}). In all
three the order is an order of
cut matrices, not of measures; only actual components of the reference are
subtracted, so the laws built from it stay nonnegative. In the middle range
$T_{\rm F}^b$ carries no middle density, the $F$ and $E$ laws having the same
middle law at class-three primes (\cref{lem:record:class-three}).
\end{remark}

\subsection{Reweighting to prescribed marginals}\label{subsec:additive-repair}

A comparison record must accept the queries, which is a condition on its
support, and it must keep the actual half marginals at $\iota$ of the record
it replaces, so that the difference annihilates the constants. A coupling with
the right support and slightly wrong marginals is corrected by the additive
reweighting $1+f(x)+g(y)$: the marginal equations are linear in $(f,g)$, and a
positive reweighting keeps the support. The next lemma bounds the centred norm
of such a reweighting; it holds for any additive perturbation with centred
$f,g$, not only for the solution of the marginal equations.

\begin{lemma}[Centred norm of an additive reweighting]\label{lem:record:additive-cancellation}
Let $\nu$ be a probability measure on $X\times Y$, finite sets, with marginals
$\nu_X,\nu_Y$ of full support, and let its centred conditional operator
$A\colon L^2_0(\nu_Y)\to L^2_0(\nu_X)$, $(Ag)(x)=\E_\nu[g(y)\mid x]$, have norm at
most $\rho$. Let $d\tilde\nu=(1+f(x)+g(y))\,d\nu\ge0$ with
$\E_{\nu_X}f=\E_{\nu_Y}g=0$, and let $1+u$, $1+v$ be the densities of its
marginals relative to $\nu_X$, $\nu_Y$. If $u\ge-U_X$, $v\ge-U_Y$,
$\|u\|_2\le e_X$, $\|v\|_2\le e_Y$, $\|f\|_\infty\le F_X$ and $\|g\|_\infty\le F_Y$,
with $U_X,U_Y<1$, then the centred conditional operator of $\tilde\nu$, relative
to its own marginals, has norm at most
\[
 \frac{\rho(1+F_X+F_Y)+e_Xe_Y}{\sqrt{(1-U_X)(1-U_Y)}} .
\]
The bounds $\|u\|_2\le e_X$, $\|v\|_2\le e_Y$ enter only through
$\|u\|_2\|v\|_2\le e_Xe_Y$, so the conclusion holds whenever $e_Xe_Y\ge\|u\|_2\|v\|_2$.
\end{lemma}

\begin{proof}
Apply \cref{lem:recentering}(a), with the references $\nu_X,\nu_Y$, to the form
$\mathfrak b(h,k)=\int hk\,d\tilde\nu$ (\cref{rem:recentering-uses}(3)): its
densities are $\varrho_L=1+u$ and $\varrho_R=1+v$, so for $h,k$ centred under the
marginals of $\tilde\nu$, with $h_0=h-\E_{\nu_X}h$ and $k_0=k-\E_{\nu_Y}k$,
\[
 \int hk\,d\tilde\nu=\int h_0k_0(1+f+g)\,d\nu-\ip u{h_0}\ip v{k_0} .
\]
Since $k_0$ is centred, $\|Ak_0\|\le\rho\|k_0\|$. Hence
\[
 \Bigl|\int h_0k_0\,d\nu\Bigr|=|\ip{h_0}{Ak_0}|\le\rho\|h_0\|\|k_0\|,\qquad
 \Bigl|\int h_0fk_0\,d\nu\Bigr|=|\ip{h_0f}{Ak_0}|\le F_X\rho\|h_0\|\|k_0\|,
\]
and the term with $g$ is bounded in the same way through $A^*$. By
Cauchy--Schwarz, $|\ip u{h_0}\ip v{k_0}|\le e_Xe_Y\|h_0\|\|k_0\|$. Finally
$1+u\ge1-U_X$ gives
$\|h\|^2_{\tilde\nu_X}\ge(1-U_X)\E_{\nu_X}h^2\ge(1-U_X)\|h_0\|^2$, and likewise for
$k$.
\end{proof}

\begin{lemma}[Additive marginal repair]\label{lem:tr:repair}
Let $\nu$, $A$ and $\rho<1$ be as in \cref{lem:record:additive-cancellation},
and let $K_\infty$ bound $A$ and $A^*$ as operators from $L^2$ to $L^\infty$,
$\|Ag\|_\infty\le K_\infty\|g\|_{L^2(\nu_Y)}$ and
$\|A^*h\|_\infty\le K_\infty\|h\|_{L^2(\nu_X)}$; one may take
$K_\infty=(\sup d\nu/d(\nu_X\otimes\nu_Y))^{1/2}$. Let $\mu_X=(1+u)\nu_X$ and
$\mu_Y=(1+v)\nu_Y$ be probability measures with $|u|\le U_X$, $|v|\le U_Y$,
$\|u\|_2\le e_X$ and $\|v\|_2\le e_Y$.
\begin{enumerate}[label=\textup{(\alph*)}]
\item The equations $f+Ag=u$, $g+A^*f=v$ have exactly one solution with
 $f\in L^2_0(\nu_X)$, $g\in L^2_0(\nu_Y)$, namely
 $f=(I-AA^*)^{-1}(u-Av)$, $g=v-A^*f$. The signed measure
 $\nu'=(1+f(x)+g(y))\,\nu$ has the marginals $\mu_X$, $\mu_Y$. The norms are
 bounded equation by equation:
 $\|f\|_2\le(e_X+\rho e_Y)/(1-\rho^2)$ and $\|g\|_2\le(e_Y+\rho e_X)/(1-\rho^2)$.
 If a group acts on $X$ and on $Y$, the diagonal action
 $(x,y)\mapsto(gx,gy)$ preserves $\nu$, and the two actions preserve
 $\mu_X$ and $\mu_Y$ respectively, then $f$, $g$ and $\nu'$ are invariant
 under these actions.
\item $\|f\|_\infty\le U_X+K_\infty\|g\|_2\le F_X:=U_X+K_\infty(e_Y+\rho e_X)/(1-\rho^2)$,
 and likewise $\|g\|_\infty\le F_Y$.
\item If $\|f\|_\infty+\|g\|_\infty<1$, for instance if $F_X+F_Y<1$, then $\nu'$ is
 a probability measure with the marginals $\mu_X,\mu_Y$, and
 $1+f+g\ge1-\|f\|_\infty-\|g\|_\infty>0$; so $\nu'$ has the support of $\nu$ and
 satisfies every constraint that holds $\nu$-almost surely.
\item \emph{(One-way repair.)} If $|u|,|v|\le U<1$, then
 $\|f\|_2,\|g\|_2\le U/(1-\rho)$ and $\|f\|_\infty,\|g\|_\infty\le
 F:=U(1+K_\infty/(1-\rho))$; if $2F<1$, (c) applies.
\item \emph{(Swap symmetry.)} If $X=Y$, $\nu$ is invariant under
 $(x,y)\mapsto(y,x)$ and $\mu_X=\mu_Y$, then $A^*=A$ and $f=g=s$, where $s$ is
 the centred solution of $(I+A)s=u$; $\nu'$ is swap invariant.
\item \emph{(Centred norm and difference.)} Under (c) and $U_X,U_Y<1$, $\nu'$
 satisfies the hypotheses of \cref{lem:record:additive-cancellation} with the
 bounds $U_X,U_Y,e_X,e_Y$ above (or any $e_X,e_Y$ with
 $e_Xe_Y\ge\|u\|_2\|v\|_2$) and any $F_X\ge\|f\|_\infty$, $F_Y\ge\|g\|_\infty$; so
 its centred conditional operator, relative to $\mu_X,\mu_Y$, has norm at most
 \[
  \rho'=\frac{\rho(1+F_X+F_Y)+e_Xe_Y}{\sqrt{(1-U_X)(1-U_Y)}};
 \]
 in the case of (d),
 \[
  \rho'\le\frac{\rho(1+2F)+U^2}{1-U}\le\frac{\rho(1+2F)+U^2+2FU}{1-U}.
 \]
 If $\pi$ is a coupling of $\mu_X$ and $\mu_Y$ whose centred conditional
 operator has norm at most $\rho(\pi)$, then $\nu'-\pi$ has zero marginals
 and, as an operator from $L^2(\mu_Y)$ to $L^2(\mu_X)$, norm at most
 $\rho'+\rho(\pi)$.
\item \emph{(Neumann series split after $N$ terms.)} For this part $\rho<1$ is not assumed. Let
 $\mathcal V_X\subseteq L^2_0(\nu_X)$ and $\mathcal V_Y\subseteq L^2_0(\nu_Y)$ be
 subspaces with $A\mathcal V_Y\subseteq\mathcal V_X$, $A^*\mathcal V_X\subseteq\mathcal V_Y$
 and $\|A|_{\mathcal V_Y}\|\le\rho_I<1$. If $u\in\mathcal V_X$, $v\in\mathcal V_Y$,
 $|u|,|v|\le U$ and $\|u-Av\|_2,\|v-A^*u\|_2\le e_I$, then
 \[
  f=\sum_{j\ge0}(AA^*)^j(u-Av)\in\mathcal V_X,\qquad
  g=\sum_{j\ge0}(A^*A)^j(v-A^*u)\in\mathcal V_Y
 \]
 solve the equations of (a), and for every integer $N\ge1$
 \[
  \|f\|_\infty,\ \|g\|_\infty\le F_N:=2NU+\frac{K_\infty\rho_I^{2N-1}e_I}{1-\rho_I^2} .
 \]
 If $2F_N<1$, the conclusions of (c) hold for this solution. If a group acts on
 $X$ and on $Y$, the diagonal action $(x,y)\mapsto(gx,gy)$ preserves
 $\nu$, and
 $\mathcal V_X$, $\mathcal V_Y$ consist of invariant functions, then $f$, $g$ and
 $\nu'$ are invariant.
\end{enumerate}
\end{lemma}

\begin{proof}
(a) $A$ maps centred functions to centred functions, since
$\E_{\nu_X}Ag=\E_{\nu_Y}g$. The positive operator $AA^*$ on $L^2_0(\nu_X)$ has norm
at most $\rho^2<1$, so $I-AA^*$ is invertible, and substitution gives
$f+Ag=(I-AA^*)f+Av=u$ and $g+A^*f=v$; the targets are centred because $\mu_X$
and $\mu_Y$ are probability measures. If $(f,g)$ solves the equations with
$u=v=0$, then $f=AA^*f$, so $f=0$ and $g=0$: the solution is unique. For
centred $f,g$ the marginal densities of $\nu'$ are
$1+f+\E_\nu[g(y)\mid x]=1+f+Ag=1+u$ and $1+v$. Taking norms in each equation
separately, $\|f\|\le e_X+\rho\|g\|$ and $\|g\|\le e_Y+\rho\|f\|$, so
$\|f\|\le e_X+\rho e_Y+\rho^2\|f\|$, which is the first bound; likewise for $g$.
(The norm $\rho$ of the off-diagonal operator on
$L^2_0(\nu_X)\oplus L^2_0(\nu_Y)$ alone would give only
$\|(f,g)\|\le\|(u,v)\|/(1-\rho)$, a loss of $\sqrt2$ in the one-way case.) A
symmetry of the data commutes with $A$ and $A^*$ and maps a centred solution to
a centred solution, which is therefore the same one.

(b) $f=u-Ag$, so $\|f\|_\infty\le\|u\|_\infty+K_\infty\|g\|_2$, and (a) bounds
$\|g\|_2$. If $k=d\nu/d(\nu_X\otimes\nu_Y)$, then
$|Ag(x)|^2\le\int k(x,y)^2\,\nu_Y(dy)\,\|g\|_2^2\le\sup k\,\|g\|_2^2$, because
$\int k(x,y)\,\nu_Y(dy)=1$.

(c) By (a) $\nu'$ has mass one and the marginals $\mu_X,\mu_Y$, and
$1+f+g\ge1-\|f\|_\infty-\|g\|_\infty>0$ on the support of $\nu$.

(d) With $U_X=U_Y=e_X=e_Y=U$, (a) gives $U(1+\rho)/(1-\rho^2)=U/(1-\rho)$ and (b)
gives $U+K_\infty U/(1-\rho)=F$.

(e) Under the swap the system for $(f,g)$ becomes the system for $(g,f)$; by
uniqueness $f=g$, and then $f+Af=u$.

(f) $\nu'\ge0$ by (c), $f,g$ are centred and $1+u,1+v$ are the densities of the
marginals of $\nu'$, with $u\ge-U_X$ and $v\ge-U_Y$;
\cref{lem:record:additive-cancellation} gives $\rho'$, and in the case of (d)
$F_X=F_Y=F$, $e_Xe_Y\le U^2$ and $\sqrt{(1-U)^2}=1-U$. Finally $\nu'$ and $\pi$
have the same marginals; their operators relative to $\mu_X,\mu_Y$ agree on the
constants, both fix them on both sides, and both map centred functions to
centred functions, so the difference is the difference of the centred parts,
of norm at most $\rho'+\rho(\pi)$.

(g) Since $A^*\mathcal V_X\subseteq\mathcal V_Y$, the restriction of $A^*$ to
$\mathcal V_X$ is the adjoint of $A|_{\mathcal V_Y}\colon\mathcal V_Y\to\mathcal V_X$;
so $AA^*$ maps $\mathcal V_X$ into itself with norm at most $\rho_I^2$, and
$A^*A$ maps $\mathcal V_Y$ into itself with norm at most $\rho_I^2$. Hence both
series converge, $f\in\mathcal V_X$, $g\in\mathcal V_Y$, $(I-AA^*)f=u-Av$, and
$A^*(AA^*)^j=(A^*A)^jA^*$ gives $g=\sum_j(A^*A)^jv-\sum_j(A^*A)^jA^*u=v-A^*f$.
Then $g+A^*f=v$ and $f+Ag=(I-AA^*)f+Av=u$. The operators $A$ and $A^*$ are
conditional expectations, hence contractions of $L^\infty$, so each of the first
$N$ terms of the series for $f$ has sup norm at most $\|u-Av\|_\infty\le2U$. For
$j\ge N$ write $(AA^*)^j(u-Av)=A\bigl[A^*(AA^*)^{j-1}(u-Av)\bigr]$; its sup norm
is at most $K_\infty\|A^*(AA^*)^{j-1}(u-Av)\|_2\le K_\infty\rho_I^{2j-1}e_I$. Summing
over $j\ge N$ gives $F_N$, and the same argument with $A$ and $A^*$ exchanged
bounds $g$. The marginal densities of $\nu'$ are $1+f+Ag=1+u$ and
$1+g+A^*f=1+v$, as in (a), and the proof of (c) applies. In the last claim
$f\in\mathcal V_X$ and $g\in\mathcal V_Y$ are invariant, and so is
$\nu'=(1+f+g)\nu$, since $\nu$ is.
\end{proof}

\begin{remark}[Blocks and the certified forms]\label{rem:repair-forms}
(1) On a block the replaced and the repaired records are tensor products over
the primes whose prime operators fix the constants of the same marginal spaces;
if every prime centred norm of both is at most $\frac12$, \cref{lem:sector}(b)
bounds the block difference by the largest prime difference: on
functions with exact support a set $S$ of $k$ primes the difference is a
sum of $k$ terms, each the prime difference at one prime of $S$ times
centred factors of norm at most $\frac12$ at the other $k-1$ primes,
hence at most $k2^{1-k}\le1$ times the largest prime difference, and on
the empty support the difference vanishes. (2) The certificates evaluate (f) in two forms. The middle merge
repairs (\cref{lem:merge-repair}) use the weaker form with $+2FU$, which the
stronger one implies, with the influence bound for $\rho(\pi)$; the one-way
repairs of the upper range (\cref{sec:upper-comparison}) use
$(\rho(1+2F)+U^2)/(1-U)$, with the smaller of the influence bound and the
spectral bound of the replaced record for $\rho(\pi)$, and in their joint-star
refinement a smaller bound for $\|u\|_2\|v\|_2$ in place of $U^2$, which (f)
allows. In every case the targets are the actual half marginals of the replaced
record, the solution of (a) keeps the support of the new graph law, and (a)
makes it invariant under the square-affine maps when the data are. (3) The
dense comparisons of the upper range and the components of its six-query
mixture (\cref{prop:dense-comparison,prop:le-mixture}) are repaired by (g) on
the functions invariant under the square-affine maps, where the centred
leading stars vanish; on these subspaces $\|A\|\le\rho_I$ with the
certified $\rho_I<1$ of the applications, and the tail of the series in
(g) beyond its first $N$ terms contributes at most
$K_\infty\rho_I^{2N-1}e_I/(1-\rho_I^2)$ to the $L^\infty$ bound $F_N$
there; the certificates record $N$ for each endpoint ($N=2$
for the dense records, $N=1$ for the components of the mixture). The changes
of these records, and those of the forest records (\cref{prop:forest-records}),
from the records they replace are bounded in the actual marginal space by
\cref{lem:record:centered-correction}: for the dense comparisons and the
mixture by its part (c), without a marginal-variance term; for the forest
records by its part (b), which also gives the full operator norms on
$L^2(m)$ entering the block-gate inequalities \eqref{eq:upper-block-gate}
for the forest records,
and for the field pairs FE and EF the symmetric part by its part (c).
\end{remark}

A repaired law is compared with the record it replaces in the space of
its actual marginals, while its bounds are proved in one common fixed
half chart. The next lemma carries out this change of space for a
reweighting $(1+f(x)+g(y))K$. It bounds separately the norm of the centred operator and
the norms of its symmetric and skew parts; and against a second law with the same
marginals, the rank-one terms produced by centring on the left and
on the right cancel, so that no marginal-variance term remains.

\begin{lemma}[Centred additive correction]\label{lem:record:centered-correction}
Let $\nu$ be a probability measure on a finite set $\mathcal H$,
$\Pi_\perp=I-\Jproj$ on $L^2(\nu)$, and $K$ a real kernel on
$\mathcal H\times\mathcal H$ with $\int K\,d(\nu\otimes\nu)=1$ (in the applications
a probability density relative to $\nu\otimes\nu$). Write $h\otimes k$ for the
kernel $h(x)k(y)$, and put $\delta_L=K\one-\one$, $\delta_R=K^*\one-\one$ and
$K_c=\Pi_\perp K\Pi_\perp$, so that
\[
 K=\one\otimes\one+\delta_L\otimes\one+\one\otimes\delta_R+K_c ,
\]
with $\delta_L$, $\delta_R$ centred. Let $\|K_c\|\le c$,
$\|\operatorname{Sym}K_c\|\le s$, $\|\operatorname{Skew}K_c\|\le a$ and
$\|\delta_L\|_\nu,\|\delta_R\|_\nu\le e$, all on $L^2(\nu)$. Let $f$, $g$ be real
functions on $\mathcal H$ with $\|f\|_\infty,\|g\|_\infty\le F$, and
$K'=(1+f(x)+g(y))K$. Here a kernel $\kappa$ denotes both the
measure $\kappa\,d(\nu\otimes\nu)$ and the operator
$h\mapsto\int\kappa(\cdot,y)h(y)\,\nu(dy)$ on $L^2(\nu)$; its half
marginals are the marginals of that measure. For a kernel $\kappa$ of
mass one with both half marginals
equal to a probability measure $m$, the \emph{centred operator} $\mathcal C_\kappa$
on $L^2(m)$ is defined by
$\ip\phi{\mathcal C_\kappa\psi}_m=\int\phi(x)\psi(y)\kappa(x,y)\,\nu(dx)\,\nu(dy)
-\E_m\phi\,\E_m\psi$; for a probability density it is the centred conditional
operator relative to the marginals.
\begin{enumerate}[label=\textup{(\alph*)}]
\item The centred correction $D=\Pi_\perp(K'-K)\Pi_\perp$ is
 \[
  D=(\Pi_\perp f)\otimes\delta_R+\delta_L\otimes(\Pi_\perp g)
  +\Pi_\perp(fK_c)\Pi_\perp+\Pi_\perp(K_cg)\Pi_\perp ,
 \]
 where $fK_c$ and $K_cg$ are the kernels $f(x)K_c(x,y)$ and $K_c(x,y)g(y)$, and
 \[
  \|D\|\le\|f\|_\nu\|\delta_R\|_\nu+\|\delta_L\|_\nu\|g\|_\nu
  +(\|f\|_\infty+\|g\|_\infty)\|K_c\|\le2F(c+e).
 \]
\item Suppose that $K'$ has mass one and both half marginals $m=(1+u_0)\nu$, with
 $1+u_0\ge e^{-\ell_m}$ and $\|u_0\|_\nu^2\le b_0$. Then
 \[
  K'-(1+u_0)\otimes(1+u_0)=K_c+D-u_0\otimes u_0 ,
 \]
 and the norms of $\mathcal C_{K'}$, of its symmetric part and of
 its skew part are at most
 \[
  e^{\ell_m}\bigl(c+2F(c+e)+b_0\bigr),\qquad
  e^{\ell_m}\bigl(s+2F(c+e)+b_0\bigr),\qquad
  e^{\ell_m}\bigl(a+2F(c+e)\bigr).
 \]
\item Suppose moreover that $O$ is a kernel of mass one with the same half
 marginals $m$ (in the applications the record that is replaced). Then
 \[
  K'-O=\Pi_\perp(K-O)\Pi_\perp+D ,
 \]
 and the symmetric and skew parts of $\mathcal C_{K'}-\mathcal C_O$ have norms at
 most $e^{\ell_m}(\|\operatorname{Sym}\Pi_\perp(K-O)\Pi_\perp\|+2F(c+e))$ and
 $e^{\ell_m}(\|\operatorname{Skew}\Pi_\perp(K-O)\Pi_\perp\|+2F(c+e))$: there is no
 marginal-variance term.
\end{enumerate}
\end{lemma}

\begin{proof}
The kernel $h(x)k(y)$ has the compression $(\Pi_\perp h)(x)(\Pi_\perp k)(y)$, so
a term that is constant in $x$ or in $y$ is killed by $\Pi_\perp\cdot\Pi_\perp$.
The four terms of the decomposition of $K$ are its four compressions by $\Jproj$
and $\Pi_\perp$: $\Jproj K\Jproj=\one\otimes\one$ because $K$ has mass one,
$\Pi_\perp K\Jproj=(K\one-\one)\otimes\one$, $\Jproj K\Pi_\perp=\one\otimes(K^*\one-\one)$,
and $\delta_L$, $\delta_R$ are centred for the same reason.

(a) Expand $(f(x)+g(y))K$ with this decomposition and compress. In
$f(x)K(x,y)$ the terms $f\otimes\one$ and $(f\delta_L)\otimes\one$ are killed, and
$f\otimes\delta_R$ and $fK_c$ give $(\Pi_\perp f)\otimes\delta_R$ and
$\Pi_\perp(fK_c)\Pi_\perp$, since $\Pi_\perp\delta_R=\delta_R$; in $g(y)K(x,y)$
the terms $\one\otimes g$ and $\one\otimes(g\delta_R)$ are killed, and
$\delta_L\otimes g$ and $K_cg$ give the other two terms. The rank-one terms have
norms $\|\Pi_\perp f\|_\nu\|\delta_R\|_\nu\le\|f\|_\nu\|\delta_R\|_\nu$ and
$\|\delta_L\|_\nu\|\Pi_\perp g\|_\nu\le\|\delta_L\|_\nu\|g\|_\nu$; the last two are
compressions of the products of $K_c$ with the multiplications by $f$ and by
$g$, of norms at most $\|f\|_\infty\|K_c\|$ and $\|g\|_\infty\|K_c\|$. With
$\|f\|_\nu\le\|f\|_\infty\le F$ and $\|g\|_\nu\le F$ this is at most $2F(c+e)$.
The bound uses only the sup norms of $f$, $g$, and no centring of them.

(b) Since $K'\one=K'^*\one=1+u_0$ and $\E_\nu u_0=0$, the kernel
$X=K'-(1+u_0)\otimes(1+u_0)$ annihilates the constants on both sides, so it equals
its compression
$\Pi_\perp K'\Pi_\perp-u_0\otimes u_0=K_c+D-u_0\otimes u_0$. By the density transport of \cref{lem:record:actual-correlation}(b), with
$\psi=1+u_0$ and $W\phi=(1+u_0)^{1/2}\phi$,
\[
 \mathcal C_{K'}=W^*M_{1+u_0}^{-1/2}XM_{1+u_0}^{-1/2}W.
\]
Here $M_{1+u_0}$ denotes multiplication by $1+u_0$. The transport multiplies
norm bounds by at most $e^{\ell_m}$ and preserves symmetric and skew parts.
The rank-one term $u_0\otimes u_0$ is self-adjoint, of norm $\|u_0\|_\nu^2\le b_0$,
so it enters the full and the symmetric norms and not the skew one; and the
symmetric and skew parts of $D$ have norms at most $\|D\|$.

(c) $K'$ and $O$ have mass one and the same half marginals, so $K'-O$
annihilates the constants on both sides and equals its compression
$\Pi_\perp(K'-O)\Pi_\perp=\Pi_\perp(K-O)\Pi_\perp+D$. As in (b),
$\mathcal C_{K'}-\mathcal C_O=W^*M_{1+u_0}^{-1/2}(K'-O)M_{1+u_0}^{-1/2}W$: the two
subtractions
$(1+u_0)\otimes(1+u_0)$ cancel. Take symmetric and skew parts.
\end{proof}

A second way to correct the marginals of a positive kernel adds instead of
reweighting. If a coupling $J$ with the prescribed (target) marginals and the
required support is at hand, a kernel with marginal defects
$\alpha,\beta$ bounded in absolute value by $M$ and satisfying
$J\beta=0$ and $J^*\alpha=0$ is corrected by the added term
$J(\alpha+\beta+2M)$ of \cref{cor:markov-completion}. This uses
neither the inversion of $I-AA^*$ nor a hypothesis $\rho<1$, and the
centred change is at most $2M(\|QRQ\|+2\|QJQ\|)$ there.

\begin{corollary}[Completion by mixing in a coupling of the target marginals]\label{cor:markov-completion}
Let $m_X$, $m_Y$ be probability measures on finite sets $X$, $Y$, and let $R\ge0$
be a density relative to $m_X\otimes m_Y$ of mass one, with marginal defects
$\alpha=\one-R\one$ and $\beta=\one-R^*\one$, $|\alpha|,|\beta|\le M$. Let $J\ge0$ be a
density relative to $m_X\otimes m_Y$ with $J\one=\one$ and $J^*\one=\one$ (a
coupling of $m_X$ and $m_Y$) such that $J\beta=0$ and $J^*\alpha=0$. Then
\[
 P=\frac{R+J\bigl(\alpha(x)+\beta(y)\bigr)+2MJ}{1+2M},
\]
where $J(\alpha(x)+\beta(y))$ is the kernel $J(x,y)(\alpha(x)+\beta(y))$, is a
probability density relative to $m_X\otimes m_Y$ with marginals $m_X$ and $m_Y$,
it vanishes wherever $R$ and $J$ both vanish, and, with $Q$ the projections off
the constants,
\[
 \|Q(P-R)Q\|\le2M\bigl(\|QRQ\|+2\|QJQ\|\bigr),
\]
the norms being those of operators from $L^2(m_Y)$ to $L^2(m_X)$. If a group acts on $X$ and on $Y$ preserving $m_X$ and
$m_Y$, and the diagonal action on $X\times Y$ preserves the densities
$R$ and $J$, then $P$ is invariant under the diagonal action.
\end{corollary}

\begin{proof}
$J(x,y)(\alpha(x)+\beta(y)+2M)\ge0$ because $\alpha+\beta\ge-2M$, so $P\ge0$. As an
operator the kernel $J(x,y)(\alpha(x)+\beta(y))$ is $M_\alpha J+JM_\beta$, with
$M_\alpha$ the multiplication by $\alpha$. By $J\one=\one$ and $J\beta=0$,
$(1+2M)P\one=(\one-\alpha)+\alpha+0+2M\one$; by $J^*\one=\one$ and $J^*\alpha=0$,
$(1+2M)P^*\one=(\one-\beta)+0+\beta+2M\one$. So $P$ has the marginals $m_X$ and
$m_Y$, and mass one; the support claim is clear. Next, $J$ maps centred functions
to centred functions, since $\ip\one{J\psi}=\ip{J^*\one}\psi$, and constants to
constants; so $JQ=QJQ$ and $QJ\Jproj=0$. Hence $QM_\alpha JQ=QM_\alpha(QJQ)$ and
$QJM_\beta Q=(QJQ)(QM_\beta Q)$, and
$\|Q(M_\alpha J+JM_\beta)Q\|\le2M\|QJQ\|$. Therefore
\[
 Q(P-R)Q=\frac{Q(M_\alpha J+JM_\beta)Q+2M\,QJQ-2M\,QRQ}{1+2M}
\]
has norm at most $2M(\|QRQ\|+2\|QJQ\|)$. Finally $\alpha$ and $\beta$ are invariant
when $R$, $m_X$ and $m_Y$ are, and the formula for $P$ is equivariant.
\end{proof}

\paragraph{Consequences and what to read next.}
From \cref{sec:positivity}: cut matrices and the order $\succeq_\vartheta$
(\cref{def:cc:order}, \eqref{eq:cc:kron}, \cref{rem:cut-symmetries}), diagonal
laws of samplewise fixed records and images under half-local maps
(\cref{lem:commoncut}(i),(v)), block brackets and their tensorization
(\cref{lem:record:block-gram}, with (d) for symmetric brackets), and the chamber
group with chamber transport (\cref{def:slot-array,lem:slot-array}). From
\cref{sec:lifts}: the sector lemma (\cref{lem:sector}), for block differences.
From \cref{sec:labellings}: the laws $\lambda_{\Blk}(G)$
(\cref{def:graphs:labeling}). From \cref{sec:cuts}: the recentering identity
(\cref{lem:recentering}(a)) and density transport
(\cref{lem:record:actual-correlation}(b)). From \cref{sec:extraction}: the joint fold and
matching rows (\cref{lem:joint-rows}(b),(c), \cref{cor:joint-row-split}).
Nothing from Part~III or Part~V is used; the references to Part~V name
applications.
